\begin{table}[H]\centering\small
\begin{tabularx}{.90\linewidth}{lX}\toprule
Symbol & Meaning in the controller and evaluation\\\midrule
$s_t$ & Controller-visible state at decision step $t$\\
$A_t$ & Candidate actions admitted by current broker authority and resource contract\\
$a_t$ & Selected proposal; execution remains subject to fresh broker admission\\
$r_t$ & Resource revision used to bind an asynchronous proposal\\
$B_t$ & Remaining time or compute budget available to a decision\\
$Y$ & Independently evaluated binary task postcondition\\
$U^P,U^F$ & Verified utility of planner and fast-policy branches\\
$C_P$ & Measured planner computation cost, including unused work where recorded\\
$\tau$ & Learned or selected delegation threshold\\
$\delta$ & Incremental utility after planner cost in a paired decision\\
\bottomrule\end{tabularx}
\caption{Notation for policy state, broker-admitted candidates, version bindings, utility, and computation cost in the evaluation.}\label{tab:notation}\end{table}

The symbols prevent internal estimates from being confused with observable task outcomes. The state $s_t$ contains only information delivered to the controller; it does not silently include hidden provider state. The admitted set $A_t$ can change when authority, resource revision, or candidate availability changes. Resource identity $r_t$ is compared for equality rather than treated as a similarity score. Success $Y$ is assigned by an independent postcondition evaluator after effects are observed. For paired branches, $U^P-U^F$ estimates local incremental value only to the extent the declared initial state and continuation policy are shared; it is not automatically a full-episode counterfactual. The episode rate is computed as $\widehat p=\sum_iY_i/N$, where $N$ remains the scheduled denominator and failed or cancelled runs are retained. Latency percentiles use controller decisions rather than episodes, and paired comparisons require identical case, variant, and seed identities. The ten repetitions in a template-shared cell do not create ten independent task families. The adequacy audit therefore reports seed-varying cells and constrains statistical interpretation. Resource revision is an identity checked for equality, while the M12 prediction residual is a separate state statistic; the two meanings must not be merged when interpreting a trace. Gate threshold $\tau$ can lie above observed probability support, producing a valid never-defer policy but no learned delegation evidence.

The episode success rate is $\widehat p=\sum_iY_i/N$, where $N$ is the scheduled denominator. Failed and cancelled episodes remain counted. Latency percentiles are computed over policy decisions and cannot be substituted for total episode wall time. Paired comparisons use the same case, variant, and seed. The sampling unit matters: ten deterministic repetitions of one template-shared cell do not create ten independent task families. The adequacy audit therefore reports the observed seed-varying cells rather than presenting all repeated episodes as independent evidence.

\section{Data and data contracts}
Neuraxis binds evidence to versioned JSON contracts rather than informal filenames. Each catalog row names a schema, producer, consumer, validator, provider, source artifact, and unavailable-data behavior. An absent trace, a malformed trace, and a valid trace with a failed outcome have different meanings. Missing scheduled evidence remains visible in the denominator; it is not replaced with presumed success or removed from an aggregate. Frozen indexes and per-episode digests support audit, while reports are recomputed from their declared sources. Table~\ref{tab:contract-families} groups the 21 contracts by the evidence boundary they govern. The individual contract subsections state the identity fields, producer and consumer path, and failure behavior that make each boundary operational.

These contracts form a dependency chain rather than 21 independent scientific findings. An episode is referenced by an index, validated against a schema, checked against lineage, and then aggregated into a report. The UI renders an evidence view only after the supporting artifacts pass their own checks. Schema validity establishes structure and required keys; producer tests show that code writes the boundary; consumer tests and provider-swap tests exercise the named reader and an alternate implementation. None proves that an outcome value is true. Scientific claims still require source digests and an independent outcome evaluator.
\begin{table}[H]\centering\small
\begin{tabularx}{.95\linewidth}{lrlX}\toprule
Evidence group & Contracts & Primary identity & Consumer-facing use\\\midrule
Product identity and release & 4 & Product version and catalog digest & Determine declared identity and release availability\\
Core experiment & 5 & Index, run, episode and lineage & Reproduce and audit the 14,400-episode core denominator\\
Transfer & 5 & Separate split and source digest & Keep semantic-transfer evidence outside the core denominator\\
Operations & 4 & Run ID, authority and input digest & Inspect bounded local operation evidence\\
Model and event audits & 3 & Checkpoint or report digest & Inspect model lineage and exploratory event-response records\\\bottomrule
\end{tabularx}
\caption{Contract groups in the versioned Neuraxis catalog and the evidence boundaries each group supplies to application views.}\label{tab:contract-families}\end{table}
\subsection{Contract boundary: product-metadata}\label{sec:contract-product_metadata}
The product-metadata boundary uses schema revision 1.3.0 and is owned by Neuraxis maintainers. Its purpose is Public product identity, license, package provenance, pinned checkpoint, manuscript records, and explicit pipeline-stage sources. The schema at \texttt{product.schema.json} governs the persisted representation; its source document is \texttt{product.json}, produced by \texttt{product.json} and consumed by \texttt{check\_\allowbreak{}contract\_\allowbreak{}catalog.py}. The executable validator is \texttt{test\_\allowbreak{}catalog.py}. Producer and consumer coverage is exercised by \texttt{test\_\allowbreak{}contract\_\allowbreak{}boundaries.py} and \texttt{test\_\allowbreak{}contract\_\allowbreak{}boundaries.py}; provider substitution is checked by \texttt{test\_\allowbreak{}contract\_\allowbreak{}boundaries.py}. These checks exercise the declared boundary but do not establish the truth of scientific outcomes.

This boundary establishes the identity a reader sees and binds that identity to the code that produced each pipeline artifact. The stage map is consequential: a release must name collection, fitting, policy optimization, gate fitting, evaluation, export, and readiness owners instead of implying that one undifferentiated pipeline performs them. Manuscript declarations and checkpoint references are part of the same identity, so stale publication metadata cannot silently describe a newer executable. The schema can verify that the fields exist and have the declared forms; the product and release gates must still compare their values with the checked-out source.

Table~\ref{tab:contract-fields-product_metadata} enumerates the schema's top-level fields, requiredness, JSON types, and explicit format, range, or enumeration constraints. Units are stated only where the schema or scientific definition supplies them; metadata and identifiers have no physical unit. A JSON type is not itself a scientific measurement. The catalog's unavailable-data rule is: missing: fail metadata validation; never infer public identity or stage ownership. For status \texttt{real}, that rule is applied at the named consumer so missing evidence remains unavailable rather than being imputed. The boundary is traceable from source artifact through producer, validator, and its actual consumer.

\small
\begin{longtable}{p{.18\linewidth}p{.16\linewidth}p{.09\linewidth}p{.43\linewidth}}
\caption{Schema fields and declared constraints for the product-metadata boundary, as defined by the product's versioned contract catalog.}\label{tab:contract-fields-product_metadata}\\
\toprule Field & JSON type & Presence & Constraint or unit\\\midrule
\endfirsthead
\caption*{product-metadata schema fields continued}\\\toprule Field & JSON type & Presence & Constraint or unit\\\midrule
\endhead
\texttt{name} & constant string & required & equals Neuraxis \\
\texttt{slug} & constant string & required & equals neuraxis \\
\texttt{version} & schema-defined & required & pattern pattern-constrained \\
\texttt{tagline} & object & required & No unit or numeric range is imposed at this field; structural checks apply. \\
\texttt{repo} & constant string & required & equals fixed value \\
\texttt{visibility} & constant string & required & equals public \\
\texttt{license} & constant string & required & equals Apache-2.0 \\
\texttt{holder} & string & required & minimum length 1 \\
\texttt{year} & integer & required & minimum 2026 \\
\texttt{method\_\allowbreak{}packages} & array of string & required & No unit or numeric range is imposed at this field; structural checks apply. \\
\texttt{method\_\allowbreak{}entrypoints} & array of string & required & minimum items 1 \\
\texttt{pipeline\_\allowbreak{}stage\_\allowbreak{}sources} & object & optional & No unit or numeric range is imposed at this field; structural checks apply. \\
\texttt{model\_\allowbreak{}artifacts} & object & optional & No unit or numeric range is imposed at this field; structural checks apply. \\
\texttt{manuscripts} & array of object & optional & No unit or numeric range is imposed at this field; structural checks apply. \\
\bottomrule
\end{longtable}
\normalsize


\subsection{Contract boundary: canonical-index}\label{sec:contract-canonical_index}
The canonical-index boundary uses schema revision 1.0.0 and is owned by Research maintainers. Its purpose is Frozen 12-method, 20-family, six-variant, ten-seed core index and lineage references. The schema at \texttt{canonical-index.schema.json} governs the persisted representation; its source document is \texttt{index.json}, produced by \texttt{run.py} and consumed by \texttt{check\_\allowbreak{}artifacts.py}. The executable validator is \texttt{test\_\allowbreak{}catalog.py}. Producer and consumer coverage is exercised by \texttt{test\_\allowbreak{}contract\_\allowbreak{}boundaries.py} and \texttt{test\_\allowbreak{}contract\_\allowbreak{}boundaries.py}; provider substitution is checked by \texttt{test\_\allowbreak{}contract\_\allowbreak{}boundaries.py}. These checks exercise the declared boundary but do not establish the truth of scientific outcomes.

The index is the closed-world inventory for the core experiment. Its method, case, episode, and aggregate lists define which records the evaluator expects; a missing cell is a coverage defect rather than an implicit failure or success. Provenance digests bind the inventory to its generated source. The artifact checker recomputes identifiers, counts, hashes, and aggregate outcomes before the browser renders a scientific summary. This contract therefore preserves the 14,400 scheduled denominator while preventing a partial directory from masquerading as a complete comparison.

Table~\ref{tab:contract-fields-canonical_index} enumerates the schema's top-level fields, requiredness, JSON types, and explicit format, range, or enumeration constraints. Units are stated only where the schema or scientific definition supplies them; metadata and identifiers have no physical unit. A JSON type is not itself a scientific measurement. The catalog's unavailable-data rule is: missing: block evidence views and fail closed. For status \texttt{real}, that rule is applied at the named consumer so missing evidence remains unavailable rather than being imputed. The boundary is traceable from source artifact through producer, validator, and its actual consumer.

\small
\begin{longtable}{p{.18\linewidth}p{.16\linewidth}p{.09\linewidth}p{.43\linewidth}}
\caption{Schema fields and declared constraints for the canonical-index boundary, as defined by the product's versioned contract catalog.}\label{tab:contract-fields-canonical_index}\\
\toprule Field & JSON type & Presence & Constraint or unit\\\midrule
\endfirsthead
\caption*{canonical-index schema fields continued}\\\toprule Field & JSON type & Presence & Constraint or unit\\\midrule
\endhead
\texttt{schema\_\allowbreak{}version} & constant int & required & equals 1 \\
\texttt{generated\_\allowbreak{}at} & string & optional & format date-time \\
\texttt{provenance} & object & required & No unit or numeric range is imposed at this field; structural checks apply. \\
\texttt{methods} & array of object & required & minimum items 12 \\
\texttt{cases} & array of object & required & minimum items 20 \\
\texttt{episodes} & array of object; episodeReference & required & minimum items 14400 \\
\texttt{aggregate\_\allowbreak{}results} & array of object & required & minimum items 12 \\
\bottomrule
\end{longtable}
\normalsize


\subsection{Contract boundary: canonical-episode}\label{sec:contract-canonical_episode}
The canonical-episode boundary uses schema revision 1.0.0 and is owned by Research maintainers. Its purpose is Executed episode decisions, outcomes, effect records, split and source lineage. The schema at \texttt{episode.schema.json} governs the persisted representation; its source document is \texttt{*.json}, produced by \texttt{external:reflexmesh==0.1.2} and consumed by \texttt{artifacts.ts}. The executable validator is \texttt{check\_\allowbreak{}artifacts.py}. Producer and consumer coverage is exercised by \texttt{test\_\allowbreak{}contract\_\allowbreak{}boundaries.py} and \texttt{test\_\allowbreak{}contract\_\allowbreak{}boundaries.py}; provider substitution is checked by \texttt{test\_\allowbreak{}contract\_\allowbreak{}boundaries.py}. These checks exercise the declared boundary but do not establish the truth of scientific outcomes.

An episode is the smallest persisted unit that connects a task identity to decisions, events, metrics, an outcome, and source lineage. The evaluator reads the final postcondition independently of the policy's chosen action. A valid episode may still fail the task; schema validity is not success. The index retains the scheduled identity when an episode is absent or malformed, and the evidence view distinguishes those cases from a completed episode whose postcondition is false. This separation is necessary for cancellation, retries, and crash recovery, where an action decision and the external effect can diverge.

Table~\ref{tab:contract-fields-canonical_episode} enumerates the schema's top-level fields, requiredness, JSON types, and explicit format, range, or enumeration constraints. Units are stated only where the schema or scientific definition supplies them; metadata and identifiers have no physical unit. A JSON type is not itself a scientific measurement. The catalog's unavailable-data rule is: missing: preserve denominator and report unavailable evidence. For status \texttt{real}, that rule is applied at the named consumer so missing evidence remains unavailable rather than being imputed. The boundary is traceable from source artifact through producer, validator, and its actual consumer.

\small
\begin{longtable}{p{.18\linewidth}p{.16\linewidth}p{.09\linewidth}p{.43\linewidth}}
\caption{Schema fields and declared constraints for the canonical-episode boundary, as defined by the product's versioned contract catalog.}\label{tab:contract-fields-canonical_episode}\\
\toprule Field & JSON type & Presence & Constraint or unit\\\midrule
\endfirsthead
\caption*{canonical-episode schema fields continued}\\\toprule Field & JSON type & Presence & Constraint or unit\\\midrule
\endhead
\texttt{schema\_\allowbreak{}version} & constant int & required & equals 1 \\
\texttt{run\_\allowbreak{}id} & string & required & minimum length 1 \\
\texttt{case\_\allowbreak{}id} & string & required & minimum length 1 \\
\texttt{method\_\allowbreak{}id} & string & required & pattern pattern-constrained \\
\texttt{variant} & string & required & minimum length 1 \\
\texttt{seed} & integer & required & No unit or numeric range is imposed at this field; structural checks apply. \\
\texttt{status} & string enum: succeeded, failed, cancelled & required & one of succeeded, failed, cancelled \\
\texttt{events} & array of schema-defined & required & minimum items 1 \\
\texttt{decisions} & array of schema-defined & required & minimum items 1 \\
\texttt{metrics} & object & required & No unit or numeric range is imposed at this field; structural checks apply. \\
\texttt{provenance} & object & required & No unit or numeric range is imposed at this field; structural checks apply. \\
\bottomrule
\end{longtable}
\normalsize


\subsection{Contract boundary: semantic-transfer}\label{sec:contract-semantic_transfer}
The semantic-transfer boundary uses schema revision 1.0.0 and is owned by Research maintainers. Its purpose is Held-out semantic transfer index, kept separate from core. The schema at \texttt{semantic-index.schema.json} governs the persisted representation; its source document is \texttt{index.json}, produced by \texttt{external:reflexmesh==0.1.2} and consumed by \texttt{Scientific\allowbreak{}Evidence.tsx}. The executable validator is \texttt{check\_\allowbreak{}semantic\_\allowbreak{}transfer.py}. Producer and consumer coverage is exercised by \texttt{test\_\allowbreak{}contract\_\allowbreak{}boundaries.py} and \texttt{test\_\allowbreak{}contract\_\allowbreak{}boundaries.py}; provider substitution is checked by \texttt{test\_\allowbreak{}contract\_\allowbreak{}boundaries.py}. These checks exercise the declared boundary but do not establish the truth of scientific outcomes.

This index defines a held-out composition lane whose operation names and task arrangements differ from the core suite. Its explicit identity and lineage keep the 1,536 transfer episodes separate from the 14,400 core episodes. The separation prevents a pooled success rate from concealing a performance drop on opaque operations. The validator checks every listed case and episode against source hashes and the lane's method matrix; an absent lane disables the transfer summary instead of borrowing core outcomes.

Table~\ref{tab:contract-fields-semantic_transfer} enumerates the schema's top-level fields, requiredness, JSON types, and explicit format, range, or enumeration constraints. Units are stated only where the schema or scientific definition supplies them; metadata and identifiers have no physical unit. A JSON type is not itself a scientific measurement. The catalog's unavailable-data rule is: missing: show transfer as unavailable and never pool it. For status \texttt{real}, that rule is applied at the named consumer so missing evidence remains unavailable rather than being imputed. The boundary is traceable from source artifact through producer, validator, and its actual consumer.

\small
\begin{longtable}{p{.18\linewidth}p{.16\linewidth}p{.09\linewidth}p{.43\linewidth}}
\caption{Schema fields and declared constraints for the semantic-transfer boundary, as defined by the product's versioned contract catalog.}\label{tab:contract-fields-semantic_transfer}\\
\toprule Field & JSON type & Presence & Constraint or unit\\\midrule
\endfirsthead
\caption*{semantic-transfer schema fields continued}\\\toprule Field & JSON type & Presence & Constraint or unit\\\midrule
\endhead
\texttt{schema\_\allowbreak{}version} & constant int & required & equals 1 \\
\texttt{provenance} & object & required & No unit or numeric range is imposed at this field; structural checks apply. \\
\texttt{methods} & array of schema-defined & required & minimum items 12 \\
\texttt{cases} & array of schema-defined & required & minimum items 8 \\
\texttt{episodes} & array of object; episodeReference & required & minimum items 1536 \\
\texttt{aggregate\_\allowbreak{}results} & array of schema-defined & required & minimum items 12 \\
\bottomrule
\end{longtable}
\normalsize


\subsection{Contract boundary: semantic-transfer-episode}\label{sec:contract-semantic_transfer_episode}
The semantic-transfer-episode boundary uses schema revision 1.0.0 and is owned by Research maintainers. Its purpose is Individually executed held-out transfer episode with provenance and observed outcomes. The schema at \texttt{episode.schema.json} governs the persisted representation; its source document is \texttt{*.json}, produced by \texttt{external:reflexmesh==0.1.2} and consumed by \texttt{artifacts.ts}. The executable validator is \texttt{check\_\allowbreak{}semantic\_\allowbreak{}transfer.py}. Producer and consumer coverage is exercised by \texttt{check\_\allowbreak{}semantic\_\allowbreak{}transfer.py} and \texttt{Scientific\allowbreak{}Evidence.tsx}; provider substitution is checked by \texttt{test\_\allowbreak{}contract\_\allowbreak{}boundaries.py}. These checks exercise the declared boundary but do not establish the truth of scientific outcomes.

Each record represents one executed held-out composition case under the transfer lane. It reuses the episode envelope while preserving the transfer case identity, variant, seed, decision trace, measured outcome, and provenance. Reusing a schema does not collapse the experiment populations: the parent index and path determine the lane, and validators reject a record whose method or case does not belong to that index. This enables matched comparisons without treating the episode as a core observation.

Table~\ref{tab:contract-fields-semantic_transfer_episode} enumerates the schema's top-level fields, requiredness, JSON types, and explicit format, range, or enumeration constraints. Units are stated only where the schema or scientific definition supplies them; metadata and identifiers have no physical unit. A JSON type is not itself a scientific measurement. The catalog's unavailable-data rule is: missing: mark the episode unavailable and preserve index denominator. For status \texttt{real}, that rule is applied at the named consumer so missing evidence remains unavailable rather than being imputed. The boundary is traceable from source artifact through producer, validator, and its actual consumer.

\small
\begin{longtable}{p{.18\linewidth}p{.16\linewidth}p{.09\linewidth}p{.43\linewidth}}
\caption{Schema fields and declared constraints for the semantic-transfer-episode boundary, as defined by the product's versioned contract catalog.}\label{tab:contract-fields-semantic_transfer_episode}\\
\toprule Field & JSON type & Presence & Constraint or unit\\\midrule
\endfirsthead
\caption*{semantic-transfer-episode schema fields continued}\\\toprule Field & JSON type & Presence & Constraint or unit\\\midrule
\endhead
\texttt{schema\_\allowbreak{}version} & constant int & required & equals 1 \\
\texttt{run\_\allowbreak{}id} & string & required & minimum length 1 \\
\texttt{case\_\allowbreak{}id} & string & required & minimum length 1 \\
\texttt{method\_\allowbreak{}id} & string & required & pattern pattern-constrained \\
\texttt{variant} & string & required & minimum length 1 \\
\texttt{seed} & integer & required & No unit or numeric range is imposed at this field; structural checks apply. \\
\texttt{status} & string enum: succeeded, failed, cancelled & required & one of succeeded, failed, cancelled \\
\texttt{events} & array of schema-defined & required & minimum items 1 \\
\texttt{decisions} & array of schema-defined & required & minimum items 1 \\
\texttt{metrics} & object & required & No unit or numeric range is imposed at this field; structural checks apply. \\
\texttt{provenance} & object & required & No unit or numeric range is imposed at this field; structural checks apply. \\
\bottomrule
\end{longtable}
\normalsize


\subsection{Contract boundary: semantic-transfer-report}\label{sec:contract-semantic_transfer_report}
The semantic-transfer-report boundary uses schema revision 1.0.0 and is owned by Research maintainers. Its purpose is Hash-bound aggregate report for the held-out semantic-transfer index. The schema at \texttt{semantic-transfer-report.schema.json} governs the persisted representation; its source document is \texttt{report.json}, produced by \texttt{check\_\allowbreak{}semantic\_\allowbreak{}transfer.py} and consumed by \texttt{Scientific\allowbreak{}Evidence.tsx}. The executable validator is \texttt{check\_\allowbreak{}semantic\_\allowbreak{}transfer.py}. Producer and consumer coverage is exercised by \texttt{check\_\allowbreak{}semantic\_\allowbreak{}transfer.py} and \texttt{Scientific\allowbreak{}Evidence.tsx}; provider substitution is checked by \texttt{test\_\allowbreak{}contract\_\allowbreak{}boundaries.py}. These checks exercise the declared boundary but do not establish the truth of scientific outcomes.

This aggregate is recomputed from the separately indexed transfer episodes and carries the index digest that defines its input population. Per-method counts, family breakdowns, and paired M12 comparisons remain tied to that digest. A report whose digest no longer matches is stale even if its JSON parses. The interface suppresses the conclusion when the report is absent or inconsistent; it never substitutes core performance. The 128 matched outcomes per method and the reported wins, ties, and losses can therefore be interpreted against a known transfer denominator.

Table~\ref{tab:contract-fields-semantic_transfer_report} enumerates the schema's top-level fields, requiredness, JSON types, and explicit format, range, or enumeration constraints. Units are stated only where the schema or scientific definition supplies them; metadata and identifiers have no physical unit. A JSON type is not itself a scientific measurement. The catalog's unavailable-data rule is: missing: suppress aggregate claims and mark transfer unavailable. For status \texttt{real}, that rule is applied at the named consumer so missing evidence remains unavailable rather than being imputed. The boundary is traceable from source artifact through producer, validator, and its actual consumer.

\small
\begin{longtable}{p{.18\linewidth}p{.16\linewidth}p{.09\linewidth}p{.43\linewidth}}
\caption{Schema fields and declared constraints for the semantic-transfer-report boundary, as defined by the product's versioned contract catalog.}\label{tab:contract-fields-semantic_transfer_report}\\
\toprule Field & JSON type & Presence & Constraint or unit\\\midrule
\endfirsthead
\caption*{semantic-transfer-report schema fields continued}\\\toprule Field & JSON type & Presence & Constraint or unit\\\midrule
\endhead
\texttt{schema\_\allowbreak{}version} & constant int & required & equals 1 \\
\texttt{scope} & string & required & minimum length 1 \\
\texttt{episodes\_\allowbreak{}checked} & integer & required & minimum 1 \\
\texttt{index\_\allowbreak{}sha256} & string & required & pattern pattern-constrained \\
\texttt{lineage\_\allowbreak{}signature} & string & required & pattern pattern-constrained \\
\texttt{methods} & array of object & required & minimum items 12 \\
\texttt{perfect\_\allowbreak{}methods} & array of string & required & No unit or numeric range is imposed at this field; structural checks apply. \\
\texttt{by\_\allowbreak{}family} & array of object & required & No unit or numeric range is imposed at this field; structural checks apply. \\
\texttt{m12\_\allowbreak{}paired} & object & required & No unit or numeric range is imposed at this field; structural checks apply. \\
\texttt{seed\_\allowbreak{}variable\_\allowbreak{}cells} & integer & required & minimum 0 \\
\bottomrule
\end{longtable}
\normalsize


\subsection{Contract boundary: semantic-transfer-lineage}\label{sec:contract-semantic_transfer_lineage}
The semantic-transfer-lineage boundary uses schema revision 1.0.0 and is owned by Research maintainers. Its purpose is Source, binary, checkpoint, package and platform lineage for transfer runs. The schema at \texttt{semantic-transfer-lineage.schema.json} governs the persisted representation; its source document is \texttt{lineage.json}, produced by \texttt{check\_\allowbreak{}semantic\_\allowbreak{}transfer.py} and consumed by \texttt{check\_\allowbreak{}semantic\_\allowbreak{}transfer.py}. The executable validator is \texttt{check\_\allowbreak{}semantic\_\allowbreak{}transfer.py}. Producer and consumer coverage is exercised by \texttt{check\_\allowbreak{}semantic\_\allowbreak{}transfer.py} and \texttt{Scientific\allowbreak{}Evidence.tsx}; provider substitution is checked by \texttt{test\_\allowbreak{}contract\_\allowbreak{}boundaries.py}. These checks exercise the declared boundary but do not establish the truth of scientific outcomes.

The lineage record identifies source code, native binary, model checkpoint, experiment specification, packages, planner configuration, operating platform, Python runtime, device, and inference thread count. Those values distinguish policy behavior from changes in the executable or compute environment. A digest mismatch means the recorded transfer result cannot be reproduced from the declared inputs, so the validator marks it unverified. Hardware and package identity do not establish scientific validity by themselves, but they are prerequisites for comparing measured latency and model outputs.

Table~\ref{tab:contract-fields-semantic_transfer_lineage} enumerates the schema's top-level fields, requiredness, JSON types, and explicit format, range, or enumeration constraints. Units are stated only where the schema or scientific definition supplies them; metadata and identifiers have no physical unit. A JSON type is not itself a scientific measurement. The catalog's unavailable-data rule is: missing: fail lineage verification; mark the run unverified; do not present it as reproducible. For status \texttt{real}, that rule is applied at the named consumer so missing evidence remains unavailable rather than being imputed. The boundary is traceable from source artifact through producer, validator, and its actual consumer.

\small
\begin{longtable}{p{.18\linewidth}p{.16\linewidth}p{.09\linewidth}p{.43\linewidth}}
\caption{Schema fields and declared constraints for the semantic-transfer-lineage boundary, as defined by the product's versioned contract catalog.}\label{tab:contract-fields-semantic_transfer_lineage}\\
\toprule Field & JSON type & Presence & Constraint or unit\\\midrule
\endfirsthead
\caption*{semantic-transfer-lineage schema fields continued}\\\toprule Field & JSON type & Presence & Constraint or unit\\\midrule
\endhead
\texttt{source\_\allowbreak{}sha256} & object; hashMap & required & No unit or numeric range is imposed at this field; structural checks apply. \\
\texttt{native\_\allowbreak{}source\_\allowbreak{}sha256} & object; hashMap & required & No unit or numeric range is imposed at this field; structural checks apply. \\
\texttt{checkpoint\_\allowbreak{}sha256} & object; hashMap & required & No unit or numeric range is imposed at this field; structural checks apply. \\
\texttt{spec\_\allowbreak{}sha256} & string & required & pattern pattern-constrained \\
\texttt{native\_\allowbreak{}binary} & object & required & No unit or numeric range is imposed at this field; structural checks apply. \\
\texttt{packages} & object & required & minimum properties 1 \\
\texttt{planner} & object & required & No unit or numeric range is imposed at this field; structural checks apply. \\
\texttt{platform} & string & required & minimum length 1 \\
\texttt{python} & string & required & minimum length 1 \\
\texttt{inference\_\allowbreak{}device} & string & required & minimum length 1 \\
\texttt{torch\_\allowbreak{}inference\_\allowbreak{}threads} & integer & optional & minimum 1 \\
\bottomrule
\end{longtable}
\normalsize


\subsection{Contract boundary: semantic-routing}\label{sec:contract-semantic_routing}
The semantic-routing boundary uses schema revision 2.0.0 and is owned by Research maintainers. Its purpose is Held-out routing predictions and comparator analysis. The schema at \texttt{semantic-routing.schema.json} governs the persisted representation; its source document is \texttt{*.json}, produced by \texttt{external:reflexmesh==0.1.2} and consumed by \texttt{Scientific\allowbreak{}Evidence.tsx}. The executable validator is \texttt{check\_\allowbreak{}semantic\_\allowbreak{}routing.py}. Producer and consumer coverage is exercised by \texttt{check\_\allowbreak{}semantic\_\allowbreak{}routing.py} and \texttt{Scientific\allowbreak{}Evidence.tsx}; provider substitution is checked by \texttt{test\_\allowbreak{}contract\_\allowbreak{}boundaries.py}. These checks exercise the declared boundary but do not establish the truth of scientific outcomes.

The persisted artifact stores held-out routing predictions and their split definition, feature description, algorithm, and evaluation status. It lets the routing analysis be recomputed without conflating training rows with the predictions used to evaluate allocation. The feature field is a description, not a numeric vector; the schema version change records that distinction. A missing routing record disables claims about learned allocation. In the measured fit, the selected threshold was 1.01 and produced never-defer behavior, so this contract preserves evidence of a fitted but non-delegating gate rather than implying useful delegation.

Table~\ref{tab:contract-fields-semantic_routing} enumerates the schema's top-level fields, requiredness, JSON types, and explicit format, range, or enumeration constraints. Units are stated only where the schema or scientific definition supplies them; metadata and identifiers have no physical unit. A JSON type is not itself a scientific measurement. The catalog's unavailable-data rule is: missing: disable routing claims and show unavailable. For status \texttt{real}, that rule is applied at the named consumer so missing evidence remains unavailable rather than being imputed. The boundary is traceable from source artifact through producer, validator, and its actual consumer.

\small
\begin{longtable}{p{.18\linewidth}p{.16\linewidth}p{.09\linewidth}p{.43\linewidth}}
\caption{Schema fields and declared constraints for the semantic-routing boundary, as defined by the product's versioned contract catalog.}\label{tab:contract-fields-semantic_routing}\\
\toprule Field & JSON type & Presence & Constraint or unit\\\midrule
\endfirsthead
\caption*{semantic-routing schema fields continued}\\\toprule Field & JSON type & Presence & Constraint or unit\\\midrule
\endhead
\texttt{schema\_\allowbreak{}version} & integer & required & minimum 1 \\
\texttt{algorithm} & string & required & minimum length 1 \\
\texttt{scope} & string & required & minimum length 1 \\
\texttt{split} & object & required & No unit or numeric range is imposed at this field; structural checks apply. \\
\texttt{features} & string & required & minimum length 1 \\
\texttt{held\_\allowbreak{}out\_\allowbreak{}predictions} & array of schema-defined & required & No unit or numeric range is imposed at this field; structural checks apply. \\
\texttt{status} & string & required & minimum length 1 \\
\bottomrule
\end{longtable}
\normalsize


\subsection{Contract boundary: operation-study-index}\label{sec:contract-operation_study_index}
The operation-study-index boundary uses schema revision 1.0.0 and is owned by Operations maintainers. Its purpose is Index of native operation runs, policy, faults, findings and latency. The schema at \texttt{operation-index.schema.json} governs the persisted representation; its source document is \texttt{index.json}, produced by \texttt{check\_\allowbreak{}operations.py} and consumed by \texttt{Scientific\allowbreak{}Evidence.tsx}. The executable validator is \texttt{check\_\allowbreak{}operations.py}. Producer and consumer coverage is exercised by \texttt{test\_\allowbreak{}contract\_\allowbreak{}boundaries.py} and \texttt{test\_\allowbreak{}contract\_\allowbreak{}boundaries.py}; provider substitution is checked by \texttt{test\_\allowbreak{}contract\_\allowbreak{}boundaries.py}. These checks exercise the declared boundary but do not establish the truth of scientific outcomes.

The index names each bounded native operation run and the engine and policy under which it was executed. Its run list is the sampling frame for the accompanying detail files; the summary must not invent runs when a detail artifact is missing. The operations page validates the index and loads only the selected detail. This creates a visible distinction between a run that is not present, one that failed, and one whose recorded checks passed. It also keeps this operational evaluation separate from replayed benchmark episodes.

Table~\ref{tab:contract-fields-operation_study_index} enumerates the schema's top-level fields, requiredness, JSON types, and explicit format, range, or enumeration constraints. Units are stated only where the schema or scientific definition supplies them; metadata and identifiers have no physical unit. A JSON type is not itself a scientific measurement. The catalog's unavailable-data rule is: missing: disable study summary; do not synthesize runs. For status \texttt{real}, that rule is applied at the named consumer so missing evidence remains unavailable rather than being imputed. The boundary is traceable from source artifact through producer, validator, and its actual consumer.

\small
\begin{longtable}{p{.18\linewidth}p{.16\linewidth}p{.09\linewidth}p{.43\linewidth}}
\caption{Schema fields and declared constraints for the operation-study-index boundary, as defined by the product's versioned contract catalog.}\label{tab:contract-fields-operation_study_index}\\
\toprule Field & JSON type & Presence & Constraint or unit\\\midrule
\endfirsthead
\caption*{operation-study-index schema fields continued}\\\toprule Field & JSON type & Presence & Constraint or unit\\\midrule
\endhead
\texttt{schema\_\allowbreak{}version} & constant int & required & equals 1 \\
\texttt{boundary} & string & required & minimum length 1 \\
\texttt{engine} & string & required & minimum length 1 \\
\texttt{policy} & string & required & minimum length 1 \\
\texttt{runs} & array of object & required & minimum items 1 \\
\bottomrule
\end{longtable}
\normalsize


\subsection{Contract boundary: operation-study-detail}\label{sec:contract-operation_study_detail}
The operation-study-detail boundary uses schema revision 2.2.0 and is owned by Operations maintainers. Its purpose is Measured operation detail including graph, effect records, findings and outputs. The schema at \texttt{operation-detail.schema.json} governs the persisted representation; its source document is \texttt{*-*.json}, produced by \texttt{check\_\allowbreak{}operations.py} and consumed by \texttt{Scientific\allowbreak{}Evidence.tsx}. The executable validator is \texttt{artifact-integrity.test.ts}. Producer and consumer coverage is exercised by \texttt{check\_\allowbreak{}operations.py} and \texttt{Scientific\allowbreak{}Evidence.tsx}; provider substitution is checked by \texttt{test\_\allowbreak{}contract\_\allowbreak{}boundaries.py}. These checks exercise the declared boundary but do not establish the truth of scientific outcomes.

A detail file contains the task scope, engine, policy, status, disposition, input identities, artifacts, events, nodes, findings, metrics, and source digest for one native run. The schema was narrowed to fields present across all persisted run families, while optional output sections remain conditional on the task. Findings and metrics are observations of the declared checks, not a general safety certification. The browser checks the detail bytes against the index before rendering it, and a missing detail affects only that run rather than invalidating unrelated records.

Table~\ref{tab:contract-fields-operation_study_detail} enumerates the schema's top-level fields, requiredness, JSON types, and explicit format, range, or enumeration constraints. Units are stated only where the schema or scientific definition supplies them; metadata and identifiers have no physical unit. A JSON type is not itself a scientific measurement. The catalog's unavailable-data rule is: missing: mark that run unavailable; preserve other run evidence. For status \texttt{real}, that rule is applied at the named consumer so missing evidence remains unavailable rather than being imputed. The boundary is traceable from source artifact through producer, validator, and its actual consumer.

\small
\begin{longtable}{p{.18\linewidth}p{.16\linewidth}p{.09\linewidth}p{.43\linewidth}}
\caption{Schema fields and declared constraints for the operation-study-detail boundary, as defined by the product's versioned contract catalog.}\label{tab:contract-fields-operation_study_detail}\\
\toprule Field & JSON type & Presence & Constraint or unit\\\midrule
\endfirsthead
\caption*{operation-study-detail schema fields continued}\\\toprule Field & JSON type & Presence & Constraint or unit\\\midrule
\endhead
\texttt{schema\_\allowbreak{}version} & integer & required & minimum 1 \\
\texttt{task} & string & required & No unit or numeric range is imposed at this field; structural checks apply. \\
\texttt{scope} & string & required & minimum length 1 \\
\texttt{engine} & string & required & minimum length 1 \\
\texttt{policy} & string & required & minimum length 1 \\
\texttt{status} & string & required & No unit or numeric range is imposed at this field; structural checks apply. \\
\texttt{disposition} & string & required & No unit or numeric range is imposed at this field; structural checks apply. \\
\texttt{inputs} & array of schema-defined & required & No unit or numeric range is imposed at this field; structural checks apply. \\
\texttt{artifacts} & object & required & No unit or numeric range is imposed at this field; structural checks apply. \\
\texttt{events} & array of schema-defined & required & No unit or numeric range is imposed at this field; structural checks apply. \\
\texttt{nodes} & array of schema-defined & required & No unit or numeric range is imposed at this field; structural checks apply. \\
\texttt{findings} & array of schema-defined & required & No unit or numeric range is imposed at this field; structural checks apply. \\
\texttt{metrics} & object & required & No unit or numeric range is imposed at this field; structural checks apply. \\
\texttt{input\_\allowbreak{}sha256} & string & required & pattern pattern-constrained \\
\bottomrule
\end{longtable}
\normalsize


\subsection{Contract boundary: operations-api}\label{sec:contract-operations_api}
The operations-api boundary uses schema revision 1.0.0 and is owned by Operations maintainers. Its purpose is Bounded local native operation service lifecycle lifecycle, cancellation, export and ownership. The schema at \texttt{operations-api.schema.json} governs the persisted representation; its source document is \texttt{operations\_\allowbreak{}api.py}; \texttt{operations.ts}; \texttt{test\_\allowbreak{}operations\_\allowbreak{}api.py}, produced by \texttt{operations\_\allowbreak{}api.py} and consumed by \texttt{Operations.tsx}. The executable validator is \texttt{test\_\allowbreak{}operations\_\allowbreak{}api.py}. Producer and consumer coverage is exercised by \texttt{test\_\allowbreak{}contract\_\allowbreak{}boundaries.py} and \texttt{test\_\allowbreak{}contract\_\allowbreak{}boundaries.py}; provider substitution is checked by \texttt{test\_\allowbreak{}contract\_\allowbreak{}boundaries.py}. These checks exercise the declared boundary but do not establish the truth of scientific outcomes.

This response envelope distinguishes accepted, running, completed, failed, cancelled, and unavailable states. A nullable run identifier and reason let the service report a rejection or unavailable provider without fabricating a run. The optional result carries the typed endpoint payload, while authentication, ownership, request bounds, and effect authorization are enforced by service logic beyond this permissive envelope. In particular, a completed HTTP response is not by itself proof that the task postcondition holds; the client must inspect the run result and its verification evidence.

Table~\ref{tab:contract-fields-operations_api} enumerates the schema's top-level fields, requiredness, JSON types, and explicit format, range, or enumeration constraints. Units are stated only where the schema or scientific definition supplies them; metadata and identifiers have no physical unit. A JSON type is not itself a scientific measurement. The catalog's unavailable-data rule is: missing: return typed unavailable response; no implicit retry/effect. For status \texttt{real}, that rule is applied at the named consumer so missing evidence remains unavailable rather than being imputed. The boundary is traceable from source artifact through producer, validator, and its actual consumer.

\small
\begin{longtable}{p{.18\linewidth}p{.16\linewidth}p{.09\linewidth}p{.43\linewidth}}
\caption{Schema fields and declared constraints for the operations-api boundary, as defined by the product's versioned contract catalog.}\label{tab:contract-fields-operations_api}\\
\toprule Field & JSON type & Presence & Constraint or unit\\\midrule
\endfirsthead
\caption*{operations-api schema fields continued}\\\toprule Field & JSON type & Presence & Constraint or unit\\\midrule
\endhead
\texttt{status} & string enum: accepted, running, completed, failed, cancelled, unavailable & required & one of accepted, running, completed, failed, cancelled, unavailable \\
\texttt{run\_\allowbreak{}id} & string or null & optional & No unit or numeric range is imposed at this field; structural checks apply. \\
\texttt{reason} & string or null & optional & No unit or numeric range is imposed at this field; structural checks apply. \\
\texttt{result} & object or array or string or number or boolean or null & optional & No unit or numeric range is imposed at this field; structural checks apply. \\
\bottomrule
\end{longtable}
\normalsize


\subsection{Contract boundary: workspace-authority}\label{sec:contract-workspace_authority}
The workspace-authority boundary uses schema revision 1.0.0 and is owned by Operations maintainers. Its purpose is Authenticated workspace registrations, previews, effects and recovery. The schema at \texttt{workspace-authority.schema.json} governs the persisted representation; its source document is \texttt{Workspace.tsx}; \texttt{test\_\allowbreak{}application.py}, produced by \texttt{main.py} and consumed by \texttt{Workspace.tsx}. The executable validator is \texttt{test\_\allowbreak{}application.py}. Producer and consumer coverage is exercised by \texttt{test\_\allowbreak{}contract\_\allowbreak{}boundaries.py} and \texttt{test\_\allowbreak{}contract\_\allowbreak{}boundaries.py}; provider substitution is checked by \texttt{test\_\allowbreak{}contract\_\allowbreak{}boundaries.py}. These checks exercise the declared boundary but do not establish the truth of scientific outcomes.

This boundary binds an authenticated workspace identity to an explicit registration and authority record, with an optional preview token and expiry. A missing registration denies effect authority; it does not select a default directory. Preview identity is time-limited and consumed by the service under the current workspace state. The schema describes the envelope, while application checks validate the permitted root, current file hashes, ownership, and single-use lifecycle. A public replay artifact cannot satisfy this contract and therefore cannot gain access to a visitor's local files.

Table~\ref{tab:contract-fields-workspace_authority} enumerates the schema's top-level fields, requiredness, JSON types, and explicit format, range, or enumeration constraints. Units are stated only where the schema or scientific definition supplies them; metadata and identifiers have no physical unit. A JSON type is not itself a scientific measurement. The catalog's unavailable-data rule is: missing: deny authority and require explicit registration. For status \texttt{real}, that rule is applied at the named consumer so missing evidence remains unavailable rather than being imputed. The boundary is traceable from source artifact through producer, validator, and its actual consumer.

\small
\begin{longtable}{p{.18\linewidth}p{.16\linewidth}p{.09\linewidth}p{.43\linewidth}}
\caption{Schema fields and declared constraints for the workspace-authority boundary, as defined by the product's versioned contract catalog.}\label{tab:contract-fields-workspace_authority}\\
\toprule Field & JSON type & Presence & Constraint or unit\\\midrule
\endfirsthead
\caption*{workspace-authority schema fields continued}\\\toprule Field & JSON type & Presence & Constraint or unit\\\midrule
\endhead
\texttt{workspace\_\allowbreak{}id} & string & required & minimum length 1 \\
\texttt{registered} & boolean & required & No unit or numeric range is imposed at this field; structural checks apply. \\
\texttt{authority} & string enum: read, write, none & required & one of read, write, none \\
\texttt{preview\_\allowbreak{}token} & string or null & optional & No unit or numeric range is imposed at this field; structural checks apply. \\
\texttt{expires\_\allowbreak{}at} & string or null & optional & format date-time \\
\bottomrule
\end{longtable}
\normalsize


\subsection{Contract boundary: product-readiness}\label{sec:contract-product_readiness}
The product-readiness boundary uses schema revision 2.0.0 and is owned by Neuraxis maintainers. Its purpose is Deterministic availability calculated from catalog and repository evidence. The schema at \texttt{readiness.schema.json} governs the persisted representation; its source document is \texttt{readiness.json}, produced by \texttt{readiness.py} and consumed by \texttt{check\_\allowbreak{}readiness.py}. The executable validator is \texttt{test\_\allowbreak{}catalog.py}. Producer and consumer coverage is exercised by \texttt{test\_\allowbreak{}contract\_\allowbreak{}boundaries.py} and \texttt{test\_\allowbreak{}contract\_\allowbreak{}boundaries.py}; provider substitution is checked by \texttt{test\_\allowbreak{}contract\_\allowbreak{}boundaries.py}. These checks exercise the declared boundary but do not establish the truth of scientific outcomes.

Readiness is derived from the contract catalog and repository paths rather than maintained as a manually optimistic status. The catalog digest lets the checker detect when the evidence inventory changed after readiness was generated. Every contract row keeps its availability state and evidence path; absent code or artifacts are reported as missing instead of receiving a default pass. The result informs release checks but does not certify scientific adequacy, user acceptance, deployment, or successful external effects.

Table~\ref{tab:contract-fields-product_readiness} enumerates the schema's top-level fields, requiredness, JSON types, and explicit format, range, or enumeration constraints. Units are stated only where the schema or scientific definition supplies them; metadata and identifiers have no physical unit. A JSON type is not itself a scientific measurement. The catalog's unavailable-data rule is: missing: fail release gate; no readiness defaults. For status \texttt{real}, that rule is applied at the named consumer so missing evidence remains unavailable rather than being imputed. The boundary is traceable from source artifact through producer, validator, and its actual consumer.

\small
\begin{longtable}{p{.18\linewidth}p{.16\linewidth}p{.09\linewidth}p{.43\linewidth}}
\caption{Schema fields and declared constraints for the product-readiness boundary, as defined by the product's versioned contract catalog.}\label{tab:contract-fields-product_readiness}\\
\toprule Field & JSON type & Presence & Constraint or unit\\\midrule
\endfirsthead
\caption*{product-readiness schema fields continued}\\\toprule Field & JSON type & Presence & Constraint or unit\\\midrule
\endhead
\texttt{schema} & constant string & required & equals neuraxis.readiness/v1 \\
\texttt{schema\_\allowbreak{}version} & constant int & required & equals 1 \\
\texttt{product\_\allowbreak{}version} & string & required & pattern pattern-constrained \\
\texttt{catalog\_\allowbreak{}sha256} & string & required & pattern pattern-constrained \\
\texttt{contracts} & array of object & required & No unit or numeric range is imposed at this field; structural checks apply. \\
\bottomrule
\end{longtable}
\normalsize


\subsection{Contract boundary: static-release}\label{sec:contract-static_release}
The static-release boundary uses schema revision 1.0.0 and is owned by Release maintainers. Its purpose is Exact static release file set, hashes, product version and source revision. The schema at \texttt{static-release.schema.json} governs the persisted representation; its source document is \texttt{package\_\allowbreak{}static.py}; \texttt{deploy\_\allowbreak{}static.py}; \texttt{test\_\allowbreak{}static\_\allowbreak{}deploy.py}, produced by \texttt{package\_\allowbreak{}static.py} and consumed by \texttt{deploy\_\allowbreak{}static.py}. The executable validator is \texttt{test\_\allowbreak{}static\_\allowbreak{}deploy.py}. Producer and consumer coverage is exercised by \texttt{test\_\allowbreak{}contract\_\allowbreak{}boundaries.py} and \texttt{test\_\allowbreak{}contract\_\allowbreak{}boundaries.py}; provider substitution is checked by \texttt{test\_\allowbreak{}contract\_\allowbreak{}boundaries.py}. These checks exercise the declared boundary but do not establish the truth of scientific outcomes.

The release manifest binds a product version and source revision to the exact set of published files and their hashes. Packaging verifies that set before deployment, and the installer checks it before atomically switching the current release. A missing or altered file blocks the new release while leaving the prior known release intact. This boundary is distinct from the research index: static packaging proves byte identity and release coherence, not that the scientific outcomes are correct.

Table~\ref{tab:contract-fields-static_release} enumerates the schema's top-level fields, requiredness, JSON types, and explicit format, range, or enumeration constraints. Units are stated only where the schema or scientific definition supplies them; metadata and identifiers have no physical unit. A JSON type is not itself a scientific measurement. The catalog's unavailable-data rule is: missing: refuse publish and preserve last known release. For status \texttt{real}, that rule is applied at the named consumer so missing evidence remains unavailable rather than being imputed. The boundary is traceable from source artifact through producer, validator, and its actual consumer.

\small
\begin{longtable}{p{.18\linewidth}p{.16\linewidth}p{.09\linewidth}p{.43\linewidth}}
\caption{Schema fields and declared constraints for the static-release boundary, as defined by the product's versioned contract catalog.}\label{tab:contract-fields-static_release}\\
\toprule Field & JSON type & Presence & Constraint or unit\\\midrule
\endfirsthead
\caption*{static-release schema fields continued}\\\toprule Field & JSON type & Presence & Constraint or unit\\\midrule
\endhead
\texttt{product} & constant string & required & equals neuraxis \\
\texttt{version} & string & required & pattern pattern-constrained \\
\texttt{source\_\allowbreak{}revision} & string & required & minimum length 7 \\
\texttt{files} & array of object & required & minimum items 1 \\
\bottomrule
\end{longtable}
\normalsize


\subsection{Contract boundary: template-adoption}\label{sec:contract-template_adoption}
The template-adoption boundary uses schema revision 1.0.0 and is owned by Neuraxis maintainers. Its purpose is Product metadata, features, architecture, notebooks, target and adversarial gates. The schema at \texttt{template-adoption.schema.json} governs the persisted representation; its source document is \texttt{.template-version}; \texttt{check\_\allowbreak{}template\_\allowbreak{}adoption.py}, produced by \texttt{.template-version} and consumed by \texttt{check\_\allowbreak{}template\_\allowbreak{}adoption.py}. The executable validator is \texttt{test\_\allowbreak{}catalog.py}. Producer and consumer coverage is exercised by \texttt{test\_\allowbreak{}contract\_\allowbreak{}boundaries.py} and \texttt{test\_\allowbreak{}contract\_\allowbreak{}boundaries.py}; provider substitution is checked by \texttt{test\_\allowbreak{}contract\_\allowbreak{}boundaries.py}. These checks exercise the declared boundary but do not establish the truth of scientific outcomes.

The persisted artifact captures the adopted template version together with required feature entries, architecture pages, notebook inventory, and deployment target. The template checker derives those values from the repository and reports missing adoption evidence. It prevents a version pin alone from being treated as complete adoption: product-specific pages and validations must also exist. Updating the template version requires re-running the installed package and content gates, because the new base may add checks that the product has not yet satisfied.

Table~\ref{tab:contract-fields-template_adoption} enumerates the schema's top-level fields, requiredness, JSON types, and explicit format, range, or enumeration constraints. Units are stated only where the schema or scientific definition supplies them; metadata and identifiers have no physical unit. A JSON type is not itself a scientific measurement. The catalog's unavailable-data rule is: missing: fail adoption gate and report exact missing evidence. For status \texttt{real}, that rule is applied at the named consumer so missing evidence remains unavailable rather than being imputed. The boundary is traceable from source artifact through producer, validator, and its actual consumer.

\small
\begin{longtable}{p{.18\linewidth}p{.16\linewidth}p{.09\linewidth}p{.43\linewidth}}
\caption{Schema fields and declared constraints for the template-adoption boundary, as defined by the product's versioned contract catalog.}\label{tab:contract-fields-template_adoption}\\
\toprule Field & JSON type & Presence & Constraint or unit\\\midrule
\endfirsthead
\caption*{template-adoption schema fields continued}\\\toprule Field & JSON type & Presence & Constraint or unit\\\midrule
\endhead
\texttt{template\_\allowbreak{}version} & string & required & pattern pattern-constrained \\
\texttt{features} & array of string & required & No unit or numeric range is imposed at this field; structural checks apply. \\
\texttt{architecture} & array of string & required & No unit or numeric range is imposed at this field; structural checks apply. \\
\texttt{deployment\_\allowbreak{}target} & string & required & minimum length 1 \\
\bottomrule
\end{longtable}
\normalsize


\subsection{Contract boundary: model-checkpoint-metadata}\label{sec:contract-model_checkpoint_metadata}
The model-checkpoint-metadata boundary uses schema revision 1.1.0 and is owned by Research maintainers. Its purpose is Method checkpoint manifests and training metadata; not the binary checkpoint payloads. The schema at \texttt{model-checkpoints.schema.json} governs the persisted representation; its source document is \texttt{*.json}, produced by \texttt{check\_\allowbreak{}artifacts.py} and consumed by \texttt{pipeline}. The executable validator is \texttt{check\_\allowbreak{}artifacts.py}. Producer and consumer coverage is exercised by \texttt{check\_\allowbreak{}artifacts.py} and \texttt{pipeline}; provider substitution is checked by \texttt{test\_\allowbreak{}contract\_\allowbreak{}boundaries.py}. These checks exercise the declared boundary but do not establish the truth of scientific outcomes.

This schema accepts several declared checkpoint and training-record shapes because linear, tree, recurrent, transition, and policy-gradient models do not share one parameter layout. The allowed variants preserve the method-specific fields needed to load or audit each model. Feature names and ordered state keys are compatibility data: a shape-compatible tensor with a different feature order is not the same model. The checker verifies referenced artifacts and lineage; missing metadata disables the affected learned method instead of falling back to an undocumented checkpoint.

Table~\ref{tab:contract-fields-model_checkpoint_metadata} enumerates the schema's top-level fields, requiredness, JSON types, and explicit format, range, or enumeration constraints. Units are stated only where the schema or scientific definition supplies them; metadata and identifiers have no physical unit. A JSON type is not itself a scientific measurement. The catalog's unavailable-data rule is: missing: disable the affected learned method and report its absent checkpoint metadata. For status \texttt{real}, that rule is applied at the named consumer so missing evidence remains unavailable rather than being imputed. The boundary is traceable from source artifact through producer, validator, and its actual consumer.

\small
\begin{longtable}{p{.18\linewidth}p{.16\linewidth}p{.09\linewidth}p{.43\linewidth}}
\caption{Schema fields and declared constraints for the model-checkpoint-metadata boundary, as defined by the product's versioned contract catalog.}\label{tab:contract-fields-model_checkpoint_metadata}\\
\toprule Field & JSON type & Presence & Constraint or unit\\\midrule
\endfirsthead
\caption*{model-checkpoint-metadata schema fields continued}\\\toprule Field & JSON type & Presence & Constraint or unit\\\midrule
\endhead
\texttt{alternative: methodBundle} & object & oneOf & Required fields: shape-specific fields \\
\texttt{alternative: linear} & object & oneOf & Required fields: schema\_version, feature\_names, coef, intercept, calibration \\
\texttt{alternative: trees} & object & oneOf & Required fields: schema\_version, feature\_names, base\_margin, trees, calibration \\
\texttt{alternative: xgboost} & object & oneOf & Required fields: learner, version \\
\texttt{alternative: recurrent} & object & oneOf & Required fields: epochs, seed, device, training\_seconds, losses, temperature, hidden\_size, feature\_names, state\_keys, onnx\_max\_abs\_error \\
\texttt{alternative: transition} & object & oneOf & Required fields: epochs, losses, device, seed, optimizer, learning\_rate, batch\_size, validation\_mse, validation\_state\_mse, validation\_reward\_mse, validation\_terminal\_mse, last\_head\_target \\
\texttt{alternative: training} & object & oneOf & Required fields: seed, device, training\_seconds \\
\texttt{alternative: metadata} & object & oneOf & Required fields: hidden\_size, temperature, feature\_names, state\_keys, onnx\_max\_abs\_error \\
\bottomrule
\end{longtable}
\normalsize


\subsection{Contract boundary: benchmark-adequacy}\label{sec:contract-benchmark_adequacy}
The benchmark-adequacy boundary uses schema revision 1.0.0 and is owned by Research maintainers. Its purpose is Hash-bound audit of saturation, seed sensitivity and per-method core outcomes. The schema at \texttt{benchmark-adequacy.schema.json} governs the persisted representation; its source document is \texttt{benchmark-adequacy.json}, produced by \texttt{check\_\allowbreak{}artifacts.py} and consumed by \texttt{Benchmark\allowbreak{}Adequacy.tsx}. The executable validator is \texttt{check\_\allowbreak{}research.py}. Producer and consumer coverage is exercised by \texttt{check\_\allowbreak{}research.py} and \texttt{Benchmark\allowbreak{}Adequacy.tsx}; provider substitution is checked by \texttt{test\_\allowbreak{}contract\_\allowbreak{}boundaries.py}. These checks exercise the declared boundary but do not establish the truth of scientific outcomes.

The adequacy record quantifies whether repeated benchmark cells provide varying outcomes and binds that calculation to the core index digest. It retains the count of scheduled episodes, the 1,440 method-case-variant cells, seed-variable cells, perfect methods, and localized M12 failures. In the frozen data, 71 cells vary by seed and five methods solve all core cases. Those values diagnose saturation and limited independent variation; they do not alter episode outcomes or establish that the benchmark represents future tasks.

Table~\ref{tab:contract-fields-benchmark_adequacy} enumerates the schema's top-level fields, requiredness, JSON types, and explicit format, range, or enumeration constraints. Units are stated only where the schema or scientific definition supplies them; metadata and identifiers have no physical unit. A JSON type is not itself a scientific measurement. The catalog's unavailable-data rule is: missing: mark benchmark-adequacy evidence unavailable. For status \texttt{real}, that rule is applied at the named consumer so missing evidence remains unavailable rather than being imputed. The boundary is traceable from source artifact through producer, validator, and its actual consumer.

\small
\begin{longtable}{p{.18\linewidth}p{.16\linewidth}p{.09\linewidth}p{.43\linewidth}}
\caption{Schema fields and declared constraints for the benchmark-adequacy boundary, as defined by the product's versioned contract catalog.}\label{tab:contract-fields-benchmark_adequacy}\\
\toprule Field & JSON type & Presence & Constraint or unit\\\midrule
\endfirsthead
\caption*{benchmark-adequacy schema fields continued}\\\toprule Field & JSON type & Presence & Constraint or unit\\\midrule
\endhead
\texttt{schema\_\allowbreak{}version} & constant int & required & equals 1 \\
\texttt{scope} & string & required & minimum length 1 \\
\texttt{episodes\_\allowbreak{}checked} & constant int & required & equals 14400 \\
\texttt{case\_\allowbreak{}variant\_\allowbreak{}method\_\allowbreak{}cells} & integer & required & minimum 1 \\
\texttt{seed\_\allowbreak{}variable\_\allowbreak{}cells} & integer & required & minimum 0 \\
\texttt{index\_\allowbreak{}sha256} & string & required & pattern pattern-constrained \\
\texttt{perfect\_\allowbreak{}methods} & array of string & required & No unit or numeric range is imposed at this field; structural checks apply. \\
\texttt{m12\_\allowbreak{}failures} & array of object & required & No unit or numeric range is imposed at this field; structural checks apply. \\
\texttt{methods} & array of object & required & minimum items 12; maximum items 12 \\
\texttt{interpretation} & string & required & minimum length 1 \\
\bottomrule
\end{longtable}
\normalsize


\subsection{Contract boundary: core-experiment-lineage}\label{sec:contract-core_experiment_lineage}
The core-experiment-lineage boundary uses schema revision 1.0.0 and is owned by Research maintainers. Its purpose is Source, binary, checkpoint, package and platform lineage for frozen core results. The schema at \texttt{core-lineage.schema.json} governs the persisted representation; its source document is \texttt{lineage.json}, produced by \texttt{check\_\allowbreak{}artifacts.py} and consumed by \texttt{Scientific\allowbreak{}Evidence.tsx}. The executable validator is \texttt{check\_\allowbreak{}artifacts.py}. Producer and consumer coverage is exercised by \texttt{check\_\allowbreak{}artifacts.py} and \texttt{Scientific\allowbreak{}Evidence.tsx}; provider substitution is checked by \texttt{test\_\allowbreak{}contract\_\allowbreak{}boundaries.py}. These checks exercise the declared boundary but do not establish the truth of scientific outcomes.

The persisted artifact identifies the code and native executable, checkpoint set, experiment specification, packages, planner, operating system, Python runtime, device, and inference thread count used for core results. It allows a reviewer to distinguish an unchanged data index from a changed implementation. The artifact checker validates the declared lineage signature against source artifacts before trusting the aggregate. A missing or mismatched lineage record disables reproducibility claims but does not delete the underlying failure episodes.

Table~\ref{tab:contract-fields-core_experiment_lineage} enumerates the schema's top-level fields, requiredness, JSON types, and explicit format, range, or enumeration constraints. Units are stated only where the schema or scientific definition supplies them; metadata and identifiers have no physical unit. A JSON type is not itself a scientific measurement. The catalog's unavailable-data rule is: missing: mark lineage unverified and disable reproducibility claims. For status \texttt{real}, that rule is applied at the named consumer so missing evidence remains unavailable rather than being imputed. The boundary is traceable from source artifact through producer, validator, and its actual consumer.

\small
\begin{longtable}{p{.18\linewidth}p{.16\linewidth}p{.09\linewidth}p{.43\linewidth}}
\caption{Schema fields and declared constraints for the core-experiment-lineage boundary, as defined by the product's versioned contract catalog.}\label{tab:contract-fields-core_experiment_lineage}\\
\toprule Field & JSON type & Presence & Constraint or unit\\\midrule
\endfirsthead
\caption*{core-experiment-lineage schema fields continued}\\\toprule Field & JSON type & Presence & Constraint or unit\\\midrule
\endhead
\texttt{source\_\allowbreak{}sha256} & object; hashMap & required & No unit or numeric range is imposed at this field; structural checks apply. \\
\texttt{native\_\allowbreak{}source\_\allowbreak{}sha256} & object; hashMap & required & No unit or numeric range is imposed at this field; structural checks apply. \\
\texttt{checkpoint\_\allowbreak{}sha256} & object & required & minimum properties 1 \\
\texttt{spec\_\allowbreak{}sha256} & string & required & pattern pattern-constrained \\
\texttt{native\_\allowbreak{}binary} & object & required & No unit or numeric range is imposed at this field; structural checks apply. \\
\texttt{packages} & object & required & minimum properties 1 \\
\texttt{planner} & object & required & No unit or numeric range is imposed at this field; structural checks apply. \\
\texttt{platform} & string & required & minimum length 1 \\
\texttt{python} & string & required & minimum length 1 \\
\texttt{inference\_\allowbreak{}device} & string & required & minimum length 1 \\
\texttt{torch\_\allowbreak{}inference\_\allowbreak{}threads} & integer & optional & minimum 1 \\
\bottomrule
\end{longtable}
\normalsize


\subsection{Contract boundary: research-report}\label{sec:contract-research_report}
The research-report boundary uses schema revision 1.0.0 and is owned by Research maintainers. Its purpose is Corrected aggregate report, paired comparisons, ablations, transfer and explicit limitations. The schema at \texttt{research-report.schema.json} governs the persisted representation; its source document is \texttt{research.json}, produced by \texttt{check\_\allowbreak{}research.py} and consumed by \texttt{Research\allowbreak{}Results.tsx}. The executable validator is \texttt{check\_\allowbreak{}research.py}. Producer and consumer coverage is exercised by \texttt{check\_\allowbreak{}research.py} and \texttt{Research\allowbreak{}Results.tsx}; provider substitution is checked by \texttt{test\_\allowbreak{}contract\_\allowbreak{}boundaries.py}. These checks exercise the declared boundary but do not establish the truth of scientific outcomes.

The aggregate report joins the core and transfer lanes with ablations, paired method differences, training details, routing, transition prediction, native parity, external comparisons, a summary conclusion, and explicit limits. Its completed count is fixed by the available source inventory and its audit digest binds the result to those inputs. The report is generated from episode-level records, not edited to improve a headline. Its limitations travel with the numbers so the M12 shortfall, saturated core suite, and separate transfer result remain visible wherever the aggregate is consumed.

Table~\ref{tab:contract-fields-research_report} enumerates the schema's top-level fields, requiredness, JSON types, and explicit format, range, or enumeration constraints. Units are stated only where the schema or scientific definition supplies them; metadata and identifiers have no physical unit. A JSON type is not itself a scientific measurement. The catalog's unavailable-data rule is: missing: hide aggregate conclusions and show evidence unavailable. For status \texttt{real}, that rule is applied at the named consumer so missing evidence remains unavailable rather than being imputed. The boundary is traceable from source artifact through producer, validator, and its actual consumer.

\small
\begin{longtable}{p{.18\linewidth}p{.16\linewidth}p{.09\linewidth}p{.43\linewidth}}
\caption{Schema fields and declared constraints for the research-report boundary, as defined by the product's versioned contract catalog.}\label{tab:contract-fields-research_report}\\
\toprule Field & JSON type & Presence & Constraint or unit\\\midrule
\endfirsthead
\caption*{research-report schema fields continued}\\\toprule Field & JSON type & Presence & Constraint or unit\\\midrule
\endhead
\texttt{schema\_\allowbreak{}version} & constant int & required & equals 1 \\
\texttt{kind} & constant string & required & equals reflexmesh-corrected-science \\
\texttt{generated\_\allowbreak{}at} & string & required & format date-time \\
\texttt{audit\_\allowbreak{}sha256} & string & required & pattern pattern-constrained \\
\texttt{completed\_\allowbreak{}episodes} & constant int & required & equals 19488 \\
\texttt{core} & object; lane & required & No unit or numeric range is imposed at this field; structural checks apply. \\
\texttt{transfer} & object; lane & required & No unit or numeric range is imposed at this field; structural checks apply. \\
\texttt{ablations} & object & required & No unit or numeric range is imposed at this field; structural checks apply. \\
\texttt{paired\_\allowbreak{}M12\_\allowbreak{}minus\_\allowbreak{}method} & object & required & minimum properties 1 \\
\texttt{training} & object & required & No unit or numeric range is imposed at this field; structural checks apply. \\
\texttt{routing} & object & required & No unit or numeric range is imposed at this field; structural checks apply. \\
\texttt{transition\_\allowbreak{}validation} & object & required & No unit or numeric range is imposed at this field; structural checks apply. \\
\texttt{native\_\allowbreak{}parity} & object & required & No unit or numeric range is imposed at this field; structural checks apply. \\
\texttt{external} & object & required & No unit or numeric range is imposed at this field; structural checks apply. \\
\texttt{ver\textbackslash\{\}-dict} & string & required & minimum length 1 \\
\texttt{limits} & array of string & required & minimum items 1 \\
\bottomrule
\end{longtable}
\normalsize


\subsection{Contract boundary: event-response-audit}\label{sec:contract-event_response_audit}
The event-response-audit boundary uses schema revision 1.0.0 and is owned by Research maintainers. Its purpose is Separate exploratory scoring of event response, invariants and final verification on frozen test episodes. The schema at \texttt{event-response-report.schema.json} governs the persisted representation; its source document is \texttt{report.json}, produced by \texttt{evaluate\_\allowbreak{}event\_\allowbreak{}responses.py} and consumed by \texttt{Event\allowbreak{}Response\allowbreak{}Evidence.tsx}. The executable validator is \texttt{check\_\allowbreak{}event\_\allowbreak{}responses.py}. Producer and consumer coverage is exercised by \texttt{evaluate\_\allowbreak{}event\_\allowbreak{}responses.py} and \texttt{Event\allowbreak{}Response\allowbreak{}Evidence.tsx}; provider substitution is checked by \texttt{test\_\allowbreak{}event\_\allowbreak{}response.py}. These checks exercise the declared boundary but do not establish the truth of scientific outcomes.

This post-hoc analysis scores response selection, event invariants, re-verification, and final completion on an existing frozen test split. The contract records its source artifact and contract definitions so that the evaluator can recompute the 4,320-episode analysis. Because the event rules were authored after the original benchmark, the result is exploratory rather than a preregistered comparison. Counts for observed events and eligible episodes are kept separate; a missing event does not count as a correct response or as a task failure by itself.

Table~\ref{tab:contract-fields-event_response_audit} enumerates the schema's top-level fields, requiredness, JSON types, and explicit format, range, or enumeration constraints. Units are stated only where the schema or scientific definition supplies them; metadata and identifiers have no physical unit. A JSON type is not itself a scientific measurement. The catalog's unavailable-data rule is: missing: hide the audit and report unavailable evidence. For status \texttt{real}, that rule is applied at the named consumer so missing evidence remains unavailable rather than being imputed. The boundary is traceable from source artifact through producer, validator, and its actual consumer.

\small
\begin{longtable}{p{.18\linewidth}p{.16\linewidth}p{.09\linewidth}p{.43\linewidth}}
\caption{Schema fields and declared constraints for the event-response-audit boundary, as defined by the product's versioned contract catalog.}\label{tab:contract-fields-event_response_audit}\\
\toprule Field & JSON type & Presence & Constraint or unit\\\midrule
\endfirsthead
\caption*{event-response-audit schema fields continued}\\\toprule Field & JSON type & Presence & Constraint or unit\\\midrule
\endhead
\texttt{schema\_\allowbreak{}version} & constant int & required & equals 1 \\
\texttt{study} & constant string & required & equals event-response-audit \\
\texttt{status} & constant string & required & equals exploratory\_posthoc \\
\texttt{generated\_\allowbreak{}at} & string & required & format date-time \\
\texttt{source} & string & required & minimum length 1 \\
\texttt{contract\_\allowbreak{}source} & string & required & minimum length 1 \\
\texttt{contracts} & array of object & required & minimum items 6; maximum items 6 \\
\texttt{dimensions} & array of string & required & minimum items 5 \\
\texttt{coverage} & object & required & No unit or numeric range is imposed at this field; structural checks apply. \\
\texttt{methods} & object & required & minimum properties 12 \\
\texttt{limitations} & array of string & required & minimum items 5 \\
\bottomrule
\end{longtable}
\normalsize


\subsection{Contract boundary: event-response-index}\label{sec:contract-event_response_index}
The event-response-index boundary uses schema revision 1.0.0 and is owned by Research maintainers. Its purpose is Source-bound identity, count and digest pointer for the exploratory event-response report. The schema at \texttt{event-response-index.schema.json} governs the persisted representation; its source document is \texttt{index.json}, produced by \texttt{evaluate\_\allowbreak{}event\_\allowbreak{}responses.py} and consumed by \texttt{check\_\allowbreak{}event\_\allowbreak{}responses.py}. The executable validator is \texttt{check\_\allowbreak{}event\_\allowbreak{}responses.py}. Producer and consumer coverage is exercised by \texttt{evaluate\_\allowbreak{}event\_\allowbreak{}responses.py} and \texttt{check\_\allowbreak{}event\_\allowbreak{}responses.py}; provider substitution is checked by \texttt{test\_\allowbreak{}event\_\allowbreak{}response.py}. These checks exercise the declared boundary but do not establish the truth of scientific outcomes.

The index binds the exploratory event-response report to its byte digest, source episode count, case list, and study identity. The checker verifies both the report hash and the frozen source count before showing any result. This prevents a recomputed report from being mistaken for the exact artifact currently referenced by the product. If the index or report is missing, the audit view becomes unavailable instead of displaying an older cached summary as current evidence.

Table~\ref{tab:contract-fields-event_response_index} enumerates the schema's top-level fields, requiredness, JSON types, and explicit format, range, or enumeration constraints. Units are stated only where the schema or scientific definition supplies them; metadata and identifiers have no physical unit. A JSON type is not itself a scientific measurement. The catalog's unavailable-data rule is: missing: hide the audit and report unavailable evidence. For status \texttt{real}, that rule is applied at the named consumer so missing evidence remains unavailable rather than being imputed. The boundary is traceable from source artifact through producer, validator, and its actual consumer.

\small
\begin{longtable}{p{.18\linewidth}p{.16\linewidth}p{.09\linewidth}p{.43\linewidth}}
\caption{Schema fields and declared constraints for the event-response-index boundary, as defined by the product's versioned contract catalog.}\label{tab:contract-fields-event_response_index}\\
\toprule Field & JSON type & Presence & Constraint or unit\\\midrule
\endfirsthead
\caption*{event-response-index schema fields continued}\\\toprule Field & JSON type & Presence & Constraint or unit\\\midrule
\endhead
\texttt{schema\_\allowbreak{}version} & constant int & required & equals 1 \\
\texttt{study} & constant string & required & equals event-response-audit \\
\texttt{report} & constant string & required & equals report.json \\
\texttt{report\_\allowbreak{}sha256} & string & required & pattern pattern-constrained \\
\texttt{source\_\allowbreak{}episode\_\allowbreak{}count} & constant int & required & equals 4320 \\
\texttt{contract\_\allowbreak{}cases} & constant list & required & equals ['C07', 'C11', 'C13', 'C15', 'C18', 'C20'] \\
\texttt{status} & constant string & required & equals exploratory\_posthoc \\
\bottomrule
\end{longtable}
\normalsize


\section{System architecture}
The runtime separates a proposal path from an effect path. A method receives an observation and candidate records and emits a candidate decision, planner request, or typed stop. The broker validates the proposal against the current candidate set and policy contract, then rechecks authority and resource revision immediately before effect execution. The provider returns an effect record or typed failure. A journal records accepted transitions for replay and crash recovery. A separate verifier evaluates the task postcondition. The policy cannot grant itself workspace authority by naming a tool, supplying arguments, or returning a high confidence score.

The application service adds ownership and resource boundaries around local operations. Workspace registration establishes which paths are in scope; previews expose intended effects; service lifecycles have bounded lifecycle, cancellation, export, and retention controls. Research artifacts use a separate path: indexes and lineage records point to immutable episode files, while the UI renders validated summaries. A replayed benchmark action therefore has no local filesystem authority, and a local file task is not a scientific episode unless the benchmark runner executes it under its protocol.

The architecture accommodates asynchronous reasoning without making it authoritative. A planner request binds goal identity, observation revision, candidates, budget, and deadline. While the request is pending, the fast policy can make progress. When a response arrives, its binding and candidate identity are checked; the broker then rechecks admission at effect time. A changed binding makes the response stale and causes discard. A digest protects the listed canonical fields from accidental mismatch, but cannot prove that all causal dependencies were included. Stale-result rejection, action admission, and postcondition verification remain separate checks. Resource versioning is a consistency mechanism, not evidence of better outcomes.

Operational guarantees are scoped to declared providers and tested contracts. They do not imply that an arbitrary remote service obeyed cancellation, that an effect record proves business completion, or that a local journal reproduces unobserved external state. Such cases need provider-specific evidence and explicit statuses such as pending, unknown, rejected, and verified. This is consistent with transaction and safe-control principles, while the concrete acceptance evidence must come from the Neuraxis integration tests and operation records \cite{lamport1998,schulman2017,suttonbarto2018,amodei2016,leike2018}.

\subsection{Runtime and action authority}
Reflexmesh represents an action as a typed candidate with a tool name, validated arguments, resource revision, and policy provenance. Policies rank or nominate candidates; the broker admits only those satisfying current contracts. Effects produce effect records and append to a journal. An effect record establishes that the runtime recorded an effect under its contract; it does not establish optimality or final task success. Policy confidence, broker admission, effect execution, and independent postcondition are distinct facts. For local file operations, registration defines the authorized root, the service validates requests against it, and export or deletion is scoped to an owned run identifier. The service and broker check different boundaries: ownership is verified at service lifecycle inspection, whereas candidate authority is rechecked immediately before an effect. A preview is an intent record, not an effect record. Unknown provider state remains unknown instead of being converted into a completion flag, and a typed timeout cannot be treated as a successful cancellation. These distinctions support accurate records during concurrent file and process operations.

\subsection{Operational workbench}
The task interface supports release verification, source auditing, and data reconciliation. Workflows parse structured inputs, check tabular shape and references, compare file digests to a supplied manifest, and verify generated artifacts. Each workflow is a registered dependency graph. Its operational controller is a fixed prerequisite policy that selects the first action admitted by the current graph. Uploaded text is parsed as data and is not executed. This task interface does not invoke the twelve research controllers or the learned routing gate. The Investigations view replays committed evidence; it does not turn a replay into fresh task execution. The workbench demonstrates bounded operation records, not learned metacontrol performance. Its file parsers, CSV shape checks, duplicate-identifier checks, foreign-key checks, and release-manifest digest comparisons are deterministic validators. They produce actionable findings about supplied artifacts but do not call an LLM or infer a research policy. An operator can inspect actual verification results without attributing a decision to a learned controller that did not make it. Successful validation means the declared checks passed for the supplied files; it does not certify the source as safe, prove that an external system accepted an upload, or establish that an operation succeeded beyond its recorded effect and postcondition.

\section{Methods}
The comparison uses twelve distinct implementations through the same broker-facing candidate interface. The set spans ordered guards, behavior-tree traversal, supervised action rankers, recurrent sequence scoring, contextual reward learning, masked policy gradients, learned transition lookahead, direct planning, fixed asynchronous delegation, and cost-sensitive deferral. A common interface controls one source of variation but does not equalize training data, labels, model size, fitting compute, provider latency, or total resource cost. Therefore the results are empirical points under a recorded protocol, not a compute-equivalent tournament. Supervised rankers target demonstrated choices, bandit updates use local reward, PPO optimizes an environment reward, and model-based lookahead depends on transition accuracy; those targets are not interchangeable.

The guarded controller scans ordered predicates over admitted actions; the behavior tree composes conditions and action leaves with success, failure, and running statuses. M03 through M08 provide the learned and model-based alternatives. M09 is a direct planner on the critical path, M10 overlaps a fixed planner request with fallback control, and M11 fits a threshold over paired incremental utility. The separate method subsections specify source, equations, assumptions, parameter values, complexity, and observed performance. Algorithm~\ref{alg:guarded} gives the guard selection procedure; Algorithm~\ref{alg:async} gives asynchronous binding and accounting; Algorithm~\ref{alg:m12} gives resource-versioned request and stale-response handling.

M12 combines recurrent action preference, an effect estimate, persistent controller state, resource-version invalidation, and planner allocation. Its prediction is that under matched compute constraints and independent postconditions it improves the success-cost frontier over strong fixed and asynchronous baselines without admitting stale or unauthorized effects. The existing evaluation does not confirm this prediction. M12's implemented form and the negative result are retained in its dedicated manuscript; this report covers the method alongside its comparators and the full contract catalog.
\subsection{Guarded finite-state controller}\label{sec:method-m01}
The guarded finite-state controller (M01) is an explicit transition policy over the observable task state and the currently registered, admissible candidates. It does not fit weights, estimate confidence, or generate action arguments. The policy reads the state flags published by the environment, forms the set of operation kinds available among admitted candidates, then applies an ordered rule list. The rule order is part of the method: stale observations are refreshed before duplicate events are acknowledged; missing dependencies are resolved before cancellation or work; cancellation requests precede ordinary workflow continuation; validation precedes acceptance or quarantine; uncertain HTTP writes are reconciled before retry; and conflicts cause rebase before compare-and-swap. For data workflows it advances extract, transform, and report according to explicit completion flags. If no supported rule has an admitted candidate, it stops rather than inventing a new operation.

Let $s_t$ be the visible state record, $A_t$ the registered candidates, and $G(a,s_t)$ the broker's current admission predicate. The operation set and M01 choice are
\begin{equation}
O_t=\{\operatorname{op}(a):a\in A_t,\ G(a,s_t)=1\},\qquad
o_t^*=\operatorname{first}_{\prec}\{o:R_o(s_t)=1\land o\in O_t\},
\end{equation}
where $R_o$ is the hand-authored guard for operation $o$ and $\prec$ is the fixed priority order. Candidate selection is then restricted to matching admitted candidates:
\begin{equation}
a_t=\arg\min_{a\in A_t}\bigl(\mathbf{1}[G(a,s_t)=0]+\mathbf{1}[\operatorname{op}(a)\ne o_t^*]\bigr),
\end{equation}
with deterministic first-candidate tie resolution in the implementation. If there is no $o_t^*$, the result is a typed stop decision with reason \texttt{no\_supported\_rule}; it is not converted to a no-op effect.

The controller's cost is linear in the number of guards and candidates for a direct scan, $O(|R|+|A_t|)$, with constant memory beyond the current observation. There is no fitting split, random seed, calibration set, or checkpoint. Its reproducible parameters are the ordered guard definitions and the operation mapping in the pinned `reflexmesh` package; changing guard priority is a method change and requires a new source digest and evaluation. This method can be preferable when process obligations are stable and a state flag has an exact operational meaning. It is brittle when the state vocabulary is incomplete, two rules conflict, or a new task needs semantic interpretation. Rule coverage should be assessed against independently authored states, including unknown flag combinations and empty candidate sets, rather than only the happy path used to author each rule.

M01 is a strong baseline on this benchmark because many task families expose exactly the flags and transitions the rules encode. It solves 1,200 of 1,200 core episodes at 0.0219 ms median and 0.0394 ms p95 decision latency in the frozen run. This result is evidence about this finite operation suite, not about general software automation. It does not prove completeness under production concurrency, establish external validity, or make the rule set a universal optimum. An independent postcondition remains necessary: a correct rule choice can still encounter a failed effect, changed resource, or unresolved external write. The acceptance test therefore checks both that every selected candidate passed broker admission and that the final goal predicate is true; admission alone is not success.



\paragraph{Additional implementation analysis.} The implementation is intentionally inspectable at the guard level. A reviewer can enumerate the rule order and identify which observable field activates each branch; no latent state or fitted checkpoint obscures this mapping. This advantage depends on observation quality. If an event is missing, delayed, or represented by a value outside the guard vocabulary, the policy can take a reasonable-looking but wrong branch or stop. Consequently, state freshness and the existence of a candidate are preconditions, not incidental metadata. A test of this controller should enumerate each rule's true and false boundary, each pair of simultaneously true guards, empty candidate lists, stale revisions, and repeated execution after an effect record. The benchmark's 100% aggregate completion only establishes that the implemented rules cover the authored transitions; it cannot stand in for those boundary tests.

The action expression above uses a deterministic tie rule over the candidate list as presented by the runtime. Candidate order is therefore part of reproducibility and must be recorded or canonicalized. The admission predicate is evaluated again at the effect boundary, so a candidate can pass policy selection and later be rejected after a revision changes. The controller should then observe the rejection and recompute; retrying the same serialized decision would violate the state-dependent guard. The implementation has no explicit estimate of expected value, risk, or planner cost. That simplicity makes its cost low and its failure modes legible, while placing the burden of coverage on manually maintained rules. A held-out family with renamed actions or partial observations would be a meaningful test of whether the baseline's performance transfers beyond the fixture vocabulary.

\subsection{Behavior-tree controller}\label{sec:method-m02}
M02 represents reactive workflow logic as a behavior tree rather than reusing M01's rule interpreter. The implementation builds a selector root from task-specific sequence nodes. A condition node returns success or failure from the current observation. An action node returns running when a matching admitted candidate exists and failure when it does not. Sequence nodes stop at the first child that is not successful; selector nodes stop at the first child that is not a failure. This status propagation distinguishes an unavailable operation from a completed goal. Each tick emits a node trace containing the node name and status, allowing a reviewer to see which branch was considered and why later siblings were skipped.

For node $n$ and observation $o$, let $\sigma_n(o)$ be its status in $\{S,F,R\}$ for success, failure, and running. An action leaf for operation $q$ has
\begin{equation}
\sigma_{q}(o)=\begin{cases}R,&\exists a\in A(o):\operatorname{op}(a)=q\land G(a,o)=1,\\F,&\text{otherwise.}\end{cases}
\end{equation}
The sequence and selector combinators use short-circuit status rules; a sequence propagates its first $F$ or $R$, while a selector propagates its first $S$ or $R$. For a sequence with children $n_1,\ldots,n_k$, success requires every child to succeed:
\begin{equation}
\sigma_{\mathrm{seq}}(o)=S\iff\bigwedge_{i=1}^{k}\sigma_{n_i}(o)=S,\qquad
\sigma_{\mathrm{sel}}(o)=F\iff\bigwedge_{i=1}^{k}\sigma_{n_i}(o)=F.
\end{equation}
The selected action is the first running leaf reached by the traversal; the broker independently checks it before execution.

The tree has explicit priority branches for stale observation refresh, duplicate acknowledgement, dependency waiting, and cancellation, followed by task subtrees for validation, testing, HTTP, restart, cleanup, extraction, compare-and-swap, jobs, reconciliation, and untrusted content. Its traversal cost is $O(N)$ in the visited nodes and it stores no learned or persistent hidden state. Unlike M01, the tree's hierarchy lets operation-specific conditions be composed and recorded as named nodes. It still uses hand-authored predicates and fixed ordering; hierarchy does not itself learn a better policy or prevent a poorly ordered selector from starving a lower-priority branch. A running leaf is a proposal for an admitted action, not proof that the environment effect finished.

The distinction between tree status and task outcome is material. `running` means that the node found a candidate requiring execution; after execution, a new observation reticks the tree. A leaf that returns failure because no matching candidate is present can cause a selector to consider another subtree, so candidate absence is not silently treated as task failure. Conversely, a sequence can remain running while an action is pending; repeated ticks must not repeat an un-idempotent effect. The runtime's intent identity and effect record provide that boundary, while the tree records the decision path. Cancellation and asynchronous completion require explicit state updates; the behavior-tree abstraction does not imply that cancellation terminates an external process.

M02 is evaluated through the same native broker and independent task postconditions as other policies. Its performance on the core suite is perfect in the archived aggregate, but the suite's deterministic state flags and known action vocabulary favor explicit reactive control. Its proper use is a compositional deterministic baseline for workflows where preconditions and recovery branches are known. Evaluation should expose rule conflicts, repeated events, unavailable leaves, livelock, and starvation, and should count effect retries separately from tree ticks. The benchmark does not establish that a behavior tree is universally more interpretable or robust than a finite-state implementation; it establishes only that this independently implemented traversal meets the current fixture contracts.



\paragraph{Additional implementation analysis.} The tree implementation separates control-flow composition from task-specific predicates. Each leaf identifies an operation kind; the candidate resolver selects a currently admitted candidate for that leaf. The tick trace is useful evidence about evaluated branches, but its presence does not establish that every semantically relevant branch was considered. A compact tree can still encode the same assumptions as a long rule list, and selector ordering can starve lower branches if an earlier branch remains perpetually running. Tests should therefore cover nonterminal results, repeated ticks, branch starvation, missing leaves, and effect record-driven transitions. A tree status also has a narrower meaning than a workflow outcome: a running action node says that work is available, not that it completed.

The method is reset between episodes and has no persistent learned parameters. This makes seed changes irrelevant to its decision function when the same observation and candidate order are supplied. The tree's effective parameters are its serialized topology, conditions, operation names, and ordering; these should be content-addressed as part of lineage. If task authors modify a branch after examining test outcomes, the test family is no longer independent. M02's strong core success therefore provides a deterministic comparator, not evidence that behavior trees are intrinsically robust. A useful stress suite would perturb event order, remove a candidate after a pending status, and inject a conflicting urgent signal while preserving the same task postcondition. Runtime authority still belongs to the broker, not the tree traversal.

\subsection{Logistic imitation policy}\label{sec:method-m03}
M03 is a fitted candidate scorer trained to rank the demonstrated action among candidates that the environment has already registered. Each candidate is represented by the fixed feature vector declared by `reflexmesh.features`; observation-level features and candidate descriptors are combined in a common feature matrix. A binary logistic model estimates demonstration preference, and an independent Platt-style logistic calibrator maps its raw margin to a probability-like score on a disjoint calibration split. The declared target is the verified demonstration action, not task completion, authorization, or future reward. The policy selects the highest-scoring currently admissible candidate and returns its action identifier, scores, target description, and a diagnostic that explicitly says the score is not a safety bound.

For feature vector $x$ and demonstration label $y\in\{0,1\}$, the fitted linear margin and probability are
\begin{equation}
z(x)=w^\top x+b,\qquad p(y=1\mid x)=\operatorname{sigmoid}(z)=\frac{1}{1+e^{-z}}.
\end{equation}
The training run used 12,576 candidate rows; logistic regression used $C=5$, the `liblinear` solver, a 600-iteration cap, CPU execution, and random seed 17. Calibration used 2,136 separate candidates, logistic calibration with $C=10$ and a 300-iteration cap. At inference the choice is
\begin{equation}
a_t=\arg\max_{a\in A_t^{\mathrm{adm}}}\widehat p_{\mathrm{cal}}(y=1\mid x(o_t,a)),\qquad
\text{stop if }A_t^{\mathrm{adm}}=\varnothing.
\end{equation}
This is a ranking decision under a fixed set of choices; the maximum calibrated demonstration score is not a calibrated probability that the resulting task will succeed.

The CPU fitting is inexpensive compared with recurrent or planner inference, and scoring $m$ candidates with $d$ features costs $O(md)$. The exported linear artifact stores ordered feature names, coefficients, intercept, and calibration coefficients. The native implementation loads that explicit schema and rejects a feature-order mismatch rather than silently applying weights to different meanings. Python and Rust parity is checked on frozen vectors, including score comparison. The model contains no episode memory: reset is a no-op, so delayed events are represented only if their information appears in the current observation features. Distribution shift in candidate descriptions, missing feature values, or correlations between demonstration style and task identity can therefore alter ranking without a warning.

The calibration report for this fit uses 2,136 rows and records a Brier score of 0.000040 and log loss of 0.000582 for the demonstration-preference target. These values are not task-success metrics and may be optimistic under the particular split construction; they do not establish that the model is calibrated on transfer families. M03 reaches 1,200/1,200 core successes in the frozen evaluation, but that outcome is compatible with the suite being solvable by fixed rules and does not show that learning contributes beyond M01. The proper ablation is a matched fixed policy with the same candidate set and postcondition oracle. Use M03 when modest-data tabular ranking is useful and the feature contract is stable; use abstention or a reviewed fallback when candidates lie outside the training support.



\paragraph{Additional implementation analysis.} The training corpus contains 12,576 candidate rows, with a separate calibration sample of 2,136 rows. The configured logistic fit uses inverse regularization $C=5$, the liblinear solver, seed 17, and a 600-iteration cap; calibration uses $C=10$ and at most 300 iterations. These settings define a candidate-ranking model for the recorded feature representation, rather than a direct policy over arbitrary natural-language instructions. Row-level randomization would leak near-duplicate candidates from the same task into both fit and evaluation, so reported generalization must be interpreted in light of the experiment's group split and authored-family overlap.

For a multicandidate decision, the scorer converts candidate features into class scores, masks candidates not admitted by the runtime, and selects the highest scoring remaining candidate. The calibrated probabilities reported for a demonstration label are not probabilities that the chosen action will succeed. The measured Brier score of 0.000040 and log loss of 0.000582 concern that label target only. A low score can coexist with poor action utility when the label reflects the authored demonstrator or when labels are nearly deterministic. Calibration should be measured on genuinely held-out task families and against outcome-based utility, with reliability diagrams and subgroup counts. Because M03 has no recurrent history, it cannot use the order of prior observations except where a feature explicitly summarizes it. This method is a compact supervised comparator; it does not estimate the causal value of requesting a planner.

\subsection{Gradient-boosted action policy}\label{sec:method-m04}
M04 is a gradient-boosted decision-tree scorer over the same candidate-feature contract as M03. It fits a binary logistic objective to the demonstrated candidate and uses the tree ensemble's raw logit as an action-ranking score. The training pipeline enables categorical features and uses histogram-based split finding. A separate calibration split fits a logistic mapping over the raw margin. At inference, the policy scores only the candidates supplied in the observation and already admitted by the runtime; it does not generate new tool arguments or override the broker. This makes the comparison with M03 interpretable as a change in scoring family under a largely shared observation and action interface, though the fitting procedures and inductive biases remain different.

For additive trees $f_j$, the ensemble margin after $T$ boosting rounds is
\begin{equation}
F_T(x)=F_0(x)+\sum_{j=1}^{T}\eta f_j(x),\qquad
\widehat p(y=1\mid x)=\operatorname{sigmoid}(F_T(x)).
\end{equation}
The archived fit uses the XGBoost `binary:logistic` objective, 160 estimators, maximum depth 5, learning rate $0.08$, histogram tree method, categorical feature support, two CPU workers, and seed 17. Calibration uses 2,136 held-out candidate rows and the same disjoint Platt-style calibration procedure as the linear scorer. The policy's choice is
\begin{equation}
a_t=\arg\max_{a\in A_t^{\mathrm{adm}}}\operatorname{Calibrate}(F_T(x(o_t,a))).
\end{equation}
The score is calibrated only with respect to the demonstration label distribution represented in the calibration split; it is not a task-success probability and not an execution permission.

Tree scoring costs approximately the sum of tree depths for each candidate, $O(mTd)$ as a conservative dense bound for $m$ candidates, $T$ trees, and feature dimension $d$. Training is more expensive than the logistic baseline but remains offline. The package exports both the original model and an explicit tree representation with base margin, split feature index, threshold, child identifiers, missing-value route, leaf values, and calibration coefficients. Native scoring follows the serialized missing-value branch; a wrong branch or feature ordering changes decisions even if the model artifact hash is correct. Validation therefore compares raw margins and calibrated probabilities on the same vectors before evaluating task outcomes.

In the frozen core, M04 solves every scheduled task and, on the structural-transfer sample, solves 24/24 versus M12's 21/24. Its success on the current data does not establish broad superiority: the core has finite authored families and the transfer sample is small. The calibration report's low aggregate Brier and log-loss values refer to teacher-action preference. The sparse intermediate-confidence bins in the report also show why overall averages can obscure calibration uncertainty. Tree ensembles can exploit categorical interaction patterns that a linear model cannot represent, but they can memorize task-family identifiers or fail on unseen tool semantics. Feature leakage, family overlap, and the held-out split must therefore be audited independently of native parity. M04 is a strong tabular comparator, not a universal task planner.



\paragraph{Additional implementation analysis.} The XGBoost configuration uses 160 trees of maximum depth five, learning rate 0.08, histogram construction, categorical support, two jobs, and seed 17. The model ranks action candidates using the same candidate interface and relies on the broker mask for admissibility. The fitted tree ensemble can capture nonlinear feature interactions without requiring a recurrent state, but its predictions remain tied to the training distribution and feature schema. The run lineage must include the exact library version, category encoding, feature order, split identity, and checkpoint digest; changing categorical handling can change decisions even if the high-level method name stays constant.

On the small structural-transfer comparison M04 succeeds on all 24 cases while the proposed controller succeeds on 21. This is descriptive evidence for that composition, not proof that boosted trees dominate planners generally. It is especially important to keep structural transfer distinct from semantic transfer: M04 records only two successes among 128 cases in the latter lane. The difference indicates that the two lanes measure different generalization demands. A meaningful follow-up should separate task templates and tool semantics before feature fitting, report performance by family, and vary environment behavior within family. Compute should include model loading, feature construction, and any calibration step. XGBoost supplies a strong learned fast-policy baseline; its performance does not test whether slow reasoning is worth its cost.

\subsection{Recurrent imitation policy}\label{sec:method-m05}
M05 adds an episode-local gated recurrent unit to candidate imitation. The network consumes a compact state-feature vector at each decision, updates a 32-dimensional hidden state, embeds each candidate through a 48-unit tanh layer, concatenates candidate embedding and recurrent context, and returns one logit per candidate through a further 48-unit tanh scorer. A softmax over the current candidate set converts logits into a normalized preference distribution. Hidden state is indexed by goal ID and cleared on reset; one goal cannot silently inherit another task's context. This model represents short event history while keeping action arguments, resource versions, capabilities, and broker authority in explicit runtime contracts.

The recurrence is
\begin{equation}
h_t=\operatorname{GRU}(x_t,h_{t-1}),\quad e_{t,a}=\tanh(W_c\phi(o_t,a)+b_c),\quad
\ell_{t,a}=w_2^\top\tanh(W_1[e_{t,a};h_t]+b_1)+b_2.
\end{equation}
The candidate preference is
\begin{equation}
p_{t,a}=\frac{\exp(\ell_{t,a}/T)}{\sum_{a'\in A_t^{\mathrm{adm}}}\exp(\ell_{t,a'}/T)},\qquad
a_t=\arg\max_{a\in A_t^{\mathrm{adm}}}p_{t,a}.
\end{equation}
Training uses sequence batches of 32, cross-entropy against the teacher action, 25 epochs, Adam at $0.003$, gradient norm clipping at 1.0, and seed 17. Temperature is selected on separate calibration sequences from 41 values on an exponential grid spanning $e^{-2}$ to $e^2$. ONNX export is checked against CPU PyTorch at candidate counts 1, 3, 6, and 7 with maximum absolute error threshold $10^{-4}$.

At each time step the hidden state is carried from the GRU output; padded candidates are masked from loss. Sequences are shuffled by episode for minibatch updates, while within-episode event order is retained. The training target remains teacher action selection rather than task success. This distinction matters because imitation can reproduce a systematic teacher mistake. Calibration on demonstrated-action labels does not correct label bias or certify the downstream postcondition. A candidate absent from training may still receive a large score because softmax normalizes over the supplied candidates.

Inference cost is linear in sequence hidden size and candidate count for fixed architecture; memory is $O(H)$ per active goal for $H=32$. Reset must be called when a goal ends, and a process restart clears hidden state unless an explicit persistence contract is added. The model does not independently validate whether a delayed event belongs to the current goal. It relies on event identity and observation revisions in its inputs. M05 reaches a perfect aggregate on the frozen core matrix, but no measured advantage over M01 follows from that result. Its role is to test whether temporal state changes candidate preference on tasks with informative history. A valid follow-up would compare cases where otherwise identical current observations have different histories and verify that recurrent updates improve postconditions without state leakage.



\paragraph{Additional implementation analysis.} The recurrent scorer encodes each candidate with a 48-unit tanh projection, updates a 32-unit gated recurrent state from the goal-local history, and scores candidates with a 48-unit network. Training uses cross-entropy for 25 epochs with Adam at learning rate 0.003, batches of 32, and gradient-norm clipping at one. Temperature is selected on a separate calibration set using a 41-point grid. Native export parity is checked against the reference network with a maximum-error tolerance of $10^{-4}$. These choices make the recurrent state an explicit method resource: episode resets must clear it, goal boundaries must select the correct history, and replay must provide the same ordered observations.

A recurrent representation can summarize sequential context that a per-step classifier cannot, but the existence of hidden state does not establish useful memory. The training objective still follows demonstrations, not verified task utility. The report should therefore distinguish sequence fit, ranking behavior, outcome success, and transfer. The parity test checks implementation agreement on a finite vector set and one checkpoint; it does not measure calibration after quantization or robustness to long histories. Runtime timing should include the state update and candidate scoring but exclude unrelated planner latency only when that distinction is explicitly reported. A stronger memory test would ablate history length, perturb event order, and use held-out families where the same local state requires different actions because of earlier events. The frozen core results do not isolate such a benefit.


\paragraph{Decision path and complexity.} The evaluation resets hidden state at episode boundaries and feeds observations in recorded event order. This makes the recurrent controller reproducible from the event stream, but exposes it to truncation and ordering assumptions. A production implementation would need explicit goal-keyed state storage and a policy for partial histories; neither follows from one aggregate score. The method cost is candidate encoding plus recurrent state update. A parity tolerance and task success are separate checks. The recurrent hidden vector is useful only if it changes choices on states whose earlier history matters, so a held-out evaluation should perturb event order and compare decisions with and without history on independently authored families.

\subsection{Contextual linear bandit}\label{sec:method-m06}
M06 implements a shared linear LinUCB action-feature model. It projects candidate features into 64 dimensions with a fixed random sign matrix whose entries are $\pm1/\sqrt{64}$. The design matrix starts at identity, the reward vector at zero, and the exploration coefficient at $\alpha=0.2$. After an observed action reward, inverse covariance and coefficient vector are updated with the Sherman--Morrison identity. At decision time each admissible candidate receives predicted linear reward plus an uncertainty bonus. The policy chooses the largest upper-confidence score with deterministic tie-breaking. It is myopic: the target is one-step reward, not total completion or long-horizon utility.

For projected feature $z_t$, inverse covariance $A_t^{-1}$, reward accumulator $b_t$, and $\theta_t=A_t^{-1}b_t$, the score is
\begin{equation}
U_t(a)=z_{t,a}^{\top}\theta_t+\alpha\sqrt{z_{t,a}^{\top}A_t^{-1}z_{t,a}},\qquad
\theta_t=A_t^{-1}b_t.
\end{equation}
For observed feature $z$ and reward $r$, the update is
\begin{equation}
u=A^{-1}z,\quad A^{-1}\leftarrow A^{-1}-\frac{uu^\top}{1+z^\top u},\quad
b\leftarrow b+rz.
\end{equation}
The archived training performs 2,096 updates from executed one-step rewards, with 64 dimensions, seed 17, and identity regularization equivalent to ridge parameter 1.0.

The logged propensity is 1.0 for the selected action because inference is deterministic. That value does not create overlap for unchosen actions; inverse-propensity estimates cannot reconstruct their outcomes from these logs. The policy does not explore randomly during held-out evaluation and receives no counterfactual labels for unexecuted actions. This limits off-policy evaluation and can make early choices self-reinforcing. Candidate projection reduces compute but can alias distinct categorical tool/action combinations.

Prediction costs $O(mdq+q^2)$ for $m$ candidates, source feature dimension $d$, and projected dimension $q=64$. The confidence bonus assumes the linear reward model and feature support are informative. Terminal task success is not the same target as immediate reward, so a high LinUCB score should not be described as likely completion. This is a controlled bandit baseline for recoverable one-step decisions; multi-step delayed effects require a sequential model or a reward definition that includes future consequences. The present suite does not establish that the bandit improves on fixed rules.



\paragraph{Additional implementation analysis.} The contextual bandit projects the 64-dimensional feature vector through a fixed random-sign projection, starts with an identity inverse covariance estimate, and uses exploration coefficient 0.2. Its fitted policy records 2,096 updates. The fit uses seed 17, while the benchmark policy selects deterministically with propensity one after training. This configuration is a myopic contextual learner: it estimates reward for the current action and does not represent how an action changes later state, how a process completes asynchronously, or how a planner might change future options. The feature projection and reward definition are therefore decisive components, not implementation details.

For a LinUCB-style score, uncertainty is based on the projected feature's leverage under the accumulated design matrix. That term encourages exploration during learning, but a deployed deterministic policy with propensity one no longer randomizes. If the benchmark's task rewards are sparse or only observed after a whole trajectory, credit assignment to intermediate selections is a challenge. The update log and reward mapping should be audited for leakage from postconditions into features, and the selection propensity must be retained if any inverse-propensity estimate is later attempted. A contextual-bandit baseline is appropriate when decisions have local reward and limited action-dependent continuation, not as a drop-in replacement for an asynchronous controller. A useful evaluation would compare online regret under randomized safe exploration and log every candidate propensity; current aggregate task success alone cannot show that the learned uncertainty bonus helps.


\paragraph{Decision path and complexity.} At inference, the candidate score combines estimated reward with a leverage-based uncertainty bonus, after which only broker-admitted actions remain eligible. The random projection reduces representation cost but can discard distinctions important to rare tool preconditions. An auditable run therefore needs the projection seed and matrix, reward definition, action propensity, and update order in its training manifest. The 2,096 updates are transitions, not necessarily 2,096 independent tasks. Since deployed selection is deterministic with propensity one, fitting-time exploration must not be attributed to benchmark-time decisions. State-update cost scales with projected dimension, while scoring scales with candidate count. This baseline tests local reward learning rather than a complete long-horizon strategy.

\subsection{Masked proximal policy optimization}\label{sec:method-m07}
M07 trains a feedforward policy with Maskable PPO in the disposable software environment. At every observation it builds a Boolean mask over the padded action space, marks only currently admissible candidates, and passes that mask to the policy distribution. The network is nonrecurrent and uses separate policy and value multilayer perceptrons with hidden widths 64 and 32. The environment's native admission and task checks remain active during training. Held-out inference orders candidates consistently, applies the same mask, requests a deterministic policy action, and maps its index back to an action ID. An empty valid-action set produces a typed stop instead of selecting padding.

For trajectory sample $(s_t,a_t)$, advantage $\hat A_t$, and probability ratio $r_t(\theta)$, PPO's clipped objective is
\begin{equation}
L^{\mathrm{clip}}(\theta)=\mathbb E_t\left[\min\left(r_t(\theta)\hat A_t,
\operatorname{clip}(r_t(\theta),1-\epsilon,1+\epsilon)\hat A_t\right)\right].
\end{equation}
The sampled distribution is restricted to current valid actions by
\begin{equation}
\pi_{\theta}^{\mathrm{mask}}(a\mid s_t)=\frac{\pi_\theta(a\mid s_t)\mathbf{1}[a\in A_t^{\mathrm{adm}}]}
{\sum_{a'}\pi_\theta(a'\mid s_t)\mathbf{1}[a'\in A_t^{\mathrm{adm}}]}.
\end{equation}
The recorded fit completes 16,384 CPU timesteps, seed 29, with rollout 256, minibatch 64, five epochs, Adam rate $5\times10^{-4}$, $\gamma=0.95$, $\lambda_{\mathrm{GAE}}=0.95$, and clip range 0.2. Entropy coefficient is zero, value coefficient 0.5, and maximum gradient norm 0.5.

PPO can optimize delayed reward in a simulator, unlike a contextual bandit's one-step update, but its evidence depends on simulator fidelity and reward shaping. Exploration occurs during training, not against an arbitrary user workspace; held-out evaluation uses a deterministic trained policy. The policy mask prevents selection of candidates marked inadmissible in the observation, but an action may become stale before execution. The native broker rechecks current authority and revisions, so the mask is a learning constraint, not a replacement for authorization.

Training required approximately 840 seconds in the archived run; this offline cost should not be confused with per-decision latency. M07 solves all core episodes in the frozen aggregate, but the saturated suite cannot isolate whether PPO learning is necessary. Evaluation should include independently varied transitions and reward-hacking checks alongside success, cost, invalid-action attempts, and effect records. The clipped update limits certain policy changes on sampled data; it does not establish regret bounds under distribution shift.



\paragraph{Additional implementation analysis.} The policy uses a masked discrete action space so inadmissible candidates are removed from the categorical distribution. The training run uses 16,384 environment steps, seed 29, CPU execution, rollout length 256, minibatches of 64, five optimization epochs, learning rate 0.0005, discount 0.95, GAE parameter 0.95, and clipping 0.2. Policy and value networks have hidden layers of 64 and 32 units. Training took 840 seconds in the recorded environment. The policy is nonrecurrent; its observation features must expose every historical variable needed to select an action. The mask must be applied identically in training and evaluation to avoid invalid-action shortcuts.

PPO optimizes a clipped policy-gradient surrogate, but that objective is not the product's task-success metric. Reward shaping, episode truncation, and the distribution of terminal outcomes determine which behavior is reinforced. To reproduce the result, the report must identify the exact reward function, normalization, wrapper order, action-to-candidate map, and checkpoint. Training compute is not included in per-decision inference latency; both should be reported when judging a system intended to operate under a resource budget. A held-out evaluation should include run-to-run training variance and family-level uncertainty rather than treating deterministic episodes from one checkpoint as independent samples. PPO is a standard sequential-control comparator here; its inclusion does not establish that reinforcement learning is inherently slower, less reliable, or less suitable than supervised methods.


\paragraph{Decision path and complexity.} At inference, the policy computes logits for candidate positions, masks disallowed actions, and samples or selects from the normalized distribution according to evaluation configuration. Reproducibility depends on candidate ordering and whether stochastic sampling remains enabled. A mask mismatch can change the effective support of the policy. Its forward cost grows with candidate feature encoding and the two-layer network; it stores no recurrent state. Reproduction requires environment wrappers, reward scaling, normalization statistics, action mapping, and checkpoint digest in addition to the training seed and step count. Policy masking is a selection constraint rather than an independent authority boundary, so the external broker remains responsible for final admission.

\subsection{Learned transition and bounded lookahead}\label{sec:method-m08}
M08 trains a compact transition model and uses bounded lookahead to rank current candidates. The network maps each candidate feature vector through 64 and 48 ReLU units to predicted next-state fields, immediate reward, and a termination estimate. Training predicts observed transition fields from executed actions. The terminal output is episode termination, which can represent success or failure; it is not a task-success probability. This semantic distinction must be preserved when reading its candidate utility.

For candidate $a$, the model returns $\widehat x'_a$, $\widehat r_a$, and $\widehat d_a$. The implementation adds one bounded expansion:
\begin{equation}
Q_1(o,a)=\widehat r_a+0.8(1-\operatorname{clip}(\widehat d_a,0,1))
\max_{a'\in A(o'_a)}\widehat r_{a'}.
\end{equation}
Its predictive training objective is a weighted sum of state, reward, and terminal-head errors:
\begin{equation}
L(\phi)=\frac{1}{N}\sum_{i=1}^N\left(\|\widehat x'_i-x'_i\|_2^2+
\lambda_r(\widehat r_i-r_i)^2+\lambda_d(\widehat d_i-d_i)^2\right).
\end{equation}
The run uses Adam at 0.003, 40 epochs, batch size 128, seed 31; horizon is 2 and expansion budget is 24. Validation MSE is 0.00711 overall, with state, reward, and terminal errors persisted separately.

Lookahead regenerates candidate scores under a hypothetical state constructed from predicted fields; it does not enumerate arbitrary plans or model every effect. If the terminal estimate is wrong, the continuation multiplier can overvalue a trajectory. If the transition model is inaccurate, search can exploit its errors. Limiting expansions bounds compute but does not prevent model exploitation. Scores are labeled utility, not calibrated probability, and predicted effects are exposed as diagnostics for the selected action.

Inference work is bounded by the configured candidate expansion loop rather than a full breadth-first tree. Predictive accuracy does not by itself bound policy regret: correlated error can preserve a good ranking, while a small error near two candidate scores can reverse it. The method is useful as a model-based control comparator when transitions repeat and alternatives have delayed consequences. To establish value, a new task set must vary transitions independently and compare closed-loop outcomes, cost, and prediction errors. The current records do not show that learned lookahead surpasses the guarded finite-state policy.


\paragraph{Additional implementation analysis.} The learned transition model predicts next-state features, reward, and termination; it does not predict verified task success as a separate target. Its network has hidden layers of 64 and 48 ReLU units, is trained with Adam at learning rate 0.003 for 40 epochs with batch size 128 and seed 31, and has validation mean-squared error 0.00711. Search depth is limited to two, with a maximum of 24 expansions. These limits bound inference but also restrict the planner's capacity to represent long action dependencies. A low average feature error does not guarantee correct ranking at decision boundaries, especially when errors concentrate on rare cancellation or external-write states.

The model-based planner rolls candidate actions forward through the learned dynamics and evaluates predicted reward. Because termination and success differ, the evaluator must not treat predicted termination as a successful goal. Model error can compound across the two-step horizon; a confidence penalty or ensemble disagreement could expose some uncertainty, but neither would repair distribution shift by itself. The transition validation split and state-feature scaling must be reported alongside the aggregate MSE. A relevant test is decision-aware accuracy: compare the model's chosen first action with a real environment rollout on held-out states, and stratify by event family. The current task suite does not show that this bounded lookahead improves on M01, and reported model error alone cannot support that claim.


\paragraph{Decision path and complexity.} The rollout procedure expands candidate actions through predicted next states, truncates search at the configured depth and expansion cap, and returns the first action from the best predicted path. This bounds computation but can discard a low-ranked branch whose model error would have made it useful. The transition model and evaluator must share feature scaling and terminal semantics. Predicted termination is not verified goal completion, so an independent evaluator checks the real postcondition after execution. Search cost scales with expanded candidates times model forward cost. A decision-level diagnostic should retain selected action, predicted state, realized state, and error by family; mean-squared validation loss alone cannot show that the model ranks first actions correctly.

\subsection{Direct deliberative planner}\label{sec:method-m09}
M09 delegates each decision to a configured language-model planner. The planner receives the current bounded observation and candidate set, returns a structured decision referring to a registered candidate, and the adapter validates that response against the observation before returning it to the runtime. The benchmark configuration identifies the planner as `qwen3.5:4b`; the method records model identity and measured call duration. Free-form prose is not an executable command. Tool selection remains limited to typed candidates, and the native broker performs the final independent authority and resource checks.

Let $\mathcal{C}_t$ be the prompt-visible candidates and $Z_t$ the planner response. The adapter parses and validates a choice using
\begin{equation}
\hat a_t=\operatorname{Parse}(Z_t),\qquad
\operatorname{Valid}(\hat a_t,o_t)=\mathbf{1}[\hat a_t\in\mathcal{C}_t]
\mathbf{1}[\operatorname{Schema}(Z_t)=1].
\end{equation}
An invalid result is rejected rather than coerced to a different candidate. Planner latency is recorded as
\begin{equation}
T_{\mathrm{plan}}=(t_{\mathrm{return}}-t_{\mathrm{request}})\times10^3\ \mathrm{ms},
\end{equation}
while episode latency and broker admission remain separate measures. The baseline is synchronous: the next decision cannot use the planner result until it returns.

M09 can use broad language context to interpret combinations not encoded in a local score, but the model may return a plausible yet wrong candidate, exceed a deadline, or misread a tool description. Its output is therefore a proposal, not a verified effect. The method has no recurrent state beyond information reintroduced in each observation and prompt. Prompt length, endpoint model, decoding configuration, timeout, and candidate serialization are part of experiment identity; changing them changes the method. The prompt template and configuration should be versioned with run lineage. Candidate-ID validation cannot independently verify that semantic reasoning was correct.

The direct planner is the slow-side comparison for fixed and learned routing. In the semantic-transfer lane it solves 38 of 128 episodes in its matched comparison with M12's 18, but this sample does not define general planner quality. Its cost includes calls that time out or return invalid choices. A complete evaluation records token or equivalent provider cost, queue delay, retries, cancellation semantics, proposal admission, and whether the returned proposal was used. The system should retain a deterministic fallback for unavailable endpoints; provider failure is not a reason to relax effect authority. This direct planner is one concrete deliberative baseline, not a universal estimate for language models.



\paragraph{Additional implementation analysis.} The direct planner receives a serialized observation, task goal, budget context, and candidate list, then returns a candidate reference. A deterministic validator rejects an answer that does not identify a currently offered candidate. This protects the interface from invented operation names but cannot establish semantic correctness. The planner can choose a valid yet harmful candidate, misunderstand a precondition, or return after its answer is useful. Endpoint model, prompt, decoding, timeout, retry rule, and candidate ordering therefore belong in lineage. For the benchmark's synchronous mode, the planner call is on the critical path; latency must include request queueing and response parsing as observed by the runner.

The recorded semantic-transfer lane gives M09 38 verified successes among 128 cases, compared with 18 for M12 and two for M04. That result is bounded to the opaque-operation transfer tasks and to the pinned planner configuration. It does not show a broad advantage over every fast controller. Provider availability, output validity, and task success are separate denominators; an unavailable provider must be recorded as such and not silently replaced by a different model. Planner cost should include tokens or equivalent provider billing when available, plus wall-clock time and failed attempts. The direct planner is a useful slow-policy comparator precisely because its success and cost are visible; it is not an authority source and cannot bypass the broker.


\paragraph{Decision path and complexity.} This synchronous call has no fast fallback on its critical path: the returned choice is available only after provider completion or timeout. Local future cancellation does not prove remote inference stopped. Every attempt, invalid answer, timeout, and unused response belongs in cost accounting. Candidate-reference validation bounds the interface but does not prove semantic correctness. A valid choice must still be checked against the current resource revision and broker authority before effect execution. The implementation is a deliberative proposer behind a deterministic safety boundary, not an autonomous executor. Provider model, prompt, decoding, timeout, retry rule, and candidate order are part of method identity and should be content-addressed with each run.

\subsection{Fixed asynchronous delegation}\label{sec:method-m10}
M10 starts a planner request at the first eligible decision and concurrently uses a guarded finite-state fallback while the request is pending. A single-worker executor bounds concurrent calls for one policy instance. The request is bound to canonical fields: goal ID, observation revision, visible state, capabilities, remaining budget and deadline, and candidate records including arguments, effects, capability requirements, resource versions, and admission flags. When a result completes, it may be used only when the current binding equals the request binding and the returned action validates against the current observation. Otherwise the response is discarded and the FSM supplies the action. Background work is collected at episode end, so completed but unused reasoning is not omitted from accounting.

For canonicalized binding record $B_t$, its digest is
\begin{equation}
\beta_t=\operatorname{SHA256}(\operatorname{JSON}_{\mathrm{sorted}}(B_t)).
\end{equation}
A completed proposal $z_t$ is eligible only if
\begin{equation}
V_t(z_t)=\mathbf{1}[\beta_t=\beta_{\mathrm{request}}]
\mathbf{1}[\operatorname{ValidateDecision}(z_t,o_t)=1].
\end{equation}
This detects a change to declared fields; it does not prove usefulness or that the binding is complete. Broker admission remains authoritative at effect time.

The lifecycle is submit, continue with fallback, receive, validate, apply or discard, and account. The first call records origin and submits a planner request, then returns the FSM decision. On a later call, a ready result is returned only when its binding remains current. Otherwise staleness is recorded and the fallback continues. Reset waits for pending work so an independent episode cannot inherit unaccounted compute. Shutdown waits for the worker and cancels futures that have not started. These are local lifecycle controls; they do not guarantee that a remote provider stopped after cancellation.

Asynchrony can reduce critical-path delay when fallback actions make useful progress, but it still consumes compute and may return too late. Reporting only fallback decision time would hide background work and unused completions. M10 records call count, completion, origin binding and an accounting flag. It is fixed delegation: the request rule does not learn if each call is worth its cost. It tests overlap and stale-result rejection as baselines before adding a learned gate. If a fast action changes the state while the request is pending, the result may be unusable; operation-specific holding and cancellation behavior therefore matter. Current core outcomes do not show M10 outperforming the strong deterministic baseline.



\paragraph{Additional implementation analysis.} M10 submits a planner request at the first eligible state and continues with a guarded finite-state fallback while a single worker processes the request. The request digest uses sorted canonical JSON over goal identity, revision, visible state, capability set, budget, deadline, and candidate records. When the result completes, the controller compares the digest with the current binding and validates the returned candidate identifier. A stale result is discarded. Completed work is accounted at reset, so a result that arrives too late still contributes to planner-use and cost records. This design separates critical-path progress from total computation, but it cannot guarantee that a remote provider has stopped when a local future is cancelled.

The binding is only as complete as its declared fields. An omitted dependency can produce a matching hash while the real world has changed; the runtime broker's final admission check is therefore still required. Asynchronous execution also creates scheduling nondeterminism: whether the response is ready at a polling boundary can change the selected action. Reproduction requires recording poll cadence, thread model, completion timestamp, timeout, and fallback transitions. A fair latency comparison reports both policy decision time and total episode wall time, including unused planner work. M10 demonstrates fixed overlap, not learned allocation. Its success on a deterministic suite cannot show whether concurrency helps when fallback and planner actions compete for mutable resources.


\paragraph{Decision path and complexity.} Asynchronous execution helps only when fallback work makes useful progress during planner delay. If a fast action changes a resource bound into the request, the planner result can become stale despite being valid for the earlier state. Canonicalization needs stable field ordering and explicit treatment of missing values; otherwise the digest is not a reliable binding check. Digest cost is linear in serialized request size, and the single-worker limit bounds local concurrency. Poll timing can change whether a response is ready at a decision boundary, so both decision latency and episode wall time matter. The controller never relaxes broker admission to consume a late answer.

\subsection{Cost-sensitive learned deferral}\label{sec:method-m11}
M11 combines a fast scorer, paired-outcome benefit classifier, gain regressor, probability calibrator, threshold, and real planner. Gate features include fast confidence, margin between leading scores, budget, normalized deadline, uncertainty and stale flags, previous action failure, and history length. Paired branches begin from the same candidate feature matrix and causal history; the pipeline checks the start signature before using a pair. It explicitly does not claim cloned hidden operating-system state. Branches execute in fresh disposable environments and share a guarded continuation. This isolates the first routing decision but is not a whole-policy causal effect.

For paired outcome $i$, let $U_i^P$ and $U_i^F$ be verified utility under planner and fast branches, and $c_i^P$ the measured planner cost. The target is
\begin{equation}
\delta_i=U_i^P-c_i^P-U_i^F,\qquad y_i=\mathbf{1}[\delta_i>0].
\end{equation}
The policy delegates only when calibrated benefit exceeds threshold and predicted gain is positive, subject to budget:
\begin{equation}
\operatorname{delegate}(x)=\mathbf{1}[\hat p(y=1\mid x)>\tau]
\mathbf{1}[\hat g(x)>0]\mathbf{1}[B_t>0.1].
\end{equation}
The fit uses logistic regression ($C=1$), ridge gain regression ($\alpha=1$), and distinct fit, calibration, and threshold-selection seed groups 50000, 60000, and 70000.

The recorded M11 fit contains 228 states and seven positive incremental-utility labels; calibration also has seven positives. The selected threshold is 1.01, above observed probability support, producing never-defer behavior. Measured calibration net incremental utility is zero. This is a negative empirical result: on these records the gate does not request a planner and has not demonstrated delegation benefit. It does not show that every task population has zero planner value. The fit is small and skewed, so it does not support a robust population risk estimate.

An unexecuted alternative has no realized reward; the pipeline does not infer its outcome from planner confidence or teacher agreement. Shared continuation reduces some post-decision differences but makes the target local to one routing decision. It may miss longer-horizon interaction and hidden-state differences after reset. A stronger evaluation needs positive and negative pairs across independently sampled families, fixed threshold selection, calibrated reliability, complete cost accounting, and an untouched final test. Core task success alone is not evidence that M11's gate improves policy outcomes when the fitted gate never delegates.



\paragraph{Additional implementation analysis.} The paired target compares verified utility from a planner branch against a fast branch and subtracts measured planner cost. This is intended to make delegation depend on incremental value rather than planner confidence alone. The feature vector includes confidence, score margin, budget, normalized deadline, uncertainty and staleness flags, prior failure, and history length. The fitting and calibration branches share a canonical start signature and causal history but do not clone hidden operating-system state. They then execute in fresh disposable environments with shared guarded continuation. The estimate is thus a local first-routing-decision contrast, not the full counterfactual value of a policy over an entire episode.

The gate uses logistic classification with $C=1$, ridge gain estimation with penalty one, and separate seed groups for fitting, calibration, and threshold choice. Each group has 228 records; the positive count is seven. The selected threshold 1.01 lies above observed support and yields no planner requests. That behavior is conservative with respect to the measured cost target, but it provides no positive evidence that the learned gate can identify valuable delegation. Calibration net incremental utility is zero on its measured sample. Before retraining, the target definition, branch reset, continuation, and positive-label provenance must be audited; adding duplicated deterministic episodes would inflate counts without adding independent cases. A valid follow-up needs independently sampled task families, planner successes and harms, and an untouched test set. The current benchmark does not establish the gate's reliability or population-level cost effectiveness.


\paragraph{Decision path and complexity.} Classifier fitting, gain regression, probability calibration, and threshold selection are separate estimation stages. Their data splits must remain grouped by task identity, or paired branches from one scenario can leak across stages. The selected threshold 1.01 is a valid output of the stated objective but lies above observed probability support, making deployment degenerate to the fast route. Reporting a gate score without its call rate would conceal this behavior. Inference cost is a feature transform plus linear model evaluation; planner cost is incurred only when a branch actually runs. Since the benchmark records no gate delegation, it does not estimate the outcome or cost of a positive learned decision.

\subsection{Resource-versioned outcome-aware metacontroller}\label{sec:method-m12}
M12 composes a recurrent fast-action scorer, learned effect model, cost-sensitive benefit gate, planner, exact prerequisite handling, asynchronous request state, and revision invalidation. The fast scorer blends recurrent preference with predicted effect utility at weight 0.15. It compares predicted state features with the next observation and updates a residual exponential moving average using 0.8 for the prior residual and 0.2 for the new error. The configured surprise threshold is 0.25. Deliberation is requested when predicted incremental benefit exceeds the selected threshold and predicted gain is positive, or when the residual is surprising, or when the observation explicitly marks insufficient evidence. These are different mechanisms and should not be collapsed into an undocumented confidence score.

Before the learned choice, M12 handles exact observation prerequisites: stale state triggers refresh, duplicate events trigger acknowledgement, and missing dependencies trigger wait. Urgent cancellation can preempt pending planner work only after these checks. This ordering was corrected after review found that an urgent branch could otherwise repeat cancellation without first refreshing or acknowledging an event. For candidate $a$, the fast score and residual are
\begin{equation}
S_t(a)=p_{\mathrm{GRU}}(a\mid o_t,h_t)+0.15U_{\mathrm{effect}}(a\mid o_t),\qquad
r_t=\operatorname{mean}|x_t-\widehat x_t|.
\end{equation}
The request predicate is
\begin{equation}
D_t=\mathbf{1}[(\hat p_t>\tau\land\hat g_t>0)\lor(r_t>0.25)\lor\operatorname{uncertain}(o_t)]
\mathbf{1}[n_t<3]\mathbf{1}[B_t>0.1].
\end{equation}
The method permits at most three planner requests per episode, expires a request after at most 10 seconds or the remaining deadline, polls no longer than every 150 ms, and uses a single worker. A proposal with a changed canonical binding is discarded; the broker rechecks admission at effect time.

Four ablations independently disable memory, effect information, deferral, or invalidation. Each records 1,160/1,200 core successes, matching full M12. The fit has the same limitation as M11: 228 states, seven positive labels, and threshold 1.01, so the learned gate never defers. The implementation includes surprise and insufficient-evidence request routes, but their presence does not show that learned allocation works. In the core evaluation M12 solves 1,160/1,200; all 40 failures are long-job cancellation. Its p50 decision latency is 18.1303 ms and p95 158.9701 ms, compared with M01 at 0.0219 and 0.0394 ms. It solves 21 of 24 structural-transfer episodes, below M04's 24. In the separate semantic-transfer lane it loses 22, wins two, and ties 104 against M09 in 128 matched cases.

The performance hypothesis is unsupported by present evidence. A new evaluation must vary tasks and transitions so planner information can change verified outcomes, include harmful and late responses, and preserve competent fixed baselines. Report success, full critical-path latency, resource cost, cancellation termination, stale-proposal attempts, invalidation, and postconditions separately. These results do not support general-intelligence or biological-equivalence claims.


\paragraph{Additional implementation analysis.} M12's state includes recurrent action history, an effect prediction residual, pending planner metadata, a resource binding, and a per-episode request count. These state variables have separate reset and invalidation conditions. Episode reset clears recurrent and residual state; a resource-revision mismatch invalidates a pending planner answer; expiry discards a response after the configured timeout or remaining deadline. The design caps requests at three per episode, uses one worker, polls at intervals no longer than 150 milliseconds, and limits the request lifetime to ten seconds. Those bounds control resource consumption but do not establish that a request will return before it matters or that a provider terminates remote work after expiry.

The effect score is a weighted addition to the recurrent candidate preference, not a calibrated probability of success. The residual is an average absolute feature prediction error; its 0.25 trigger is fixed in the evaluated configuration. A high residual can signal a surprising observation, but it can also reflect expected stochasticity or poor feature scaling. The gate trigger and explicit uncertainty flag can request planning independently, so an evaluation must report which route caused each call. The observed allocation fit never delegates at its selected threshold; therefore the system's planner requests arise from the residual or explicit-uncertainty branches, not from a demonstrated learned benefit estimate. The request count and planner use should be reported alongside outcome changes.

The cancellation precondition order is operationally consequential: refresh stale state, acknowledge duplicate events, resolve missing dependencies, then choose an urgent cancellation candidate. This avoids acting on an obsolete event, but the final external termination still needs a provider effect record and independent process-state check. M12 currently selects cancellation in all 60 audited C07 traces yet verifies only 20 successes; action choice is not execution. Because component ablations retain 1,160/1,200 core successes, the core evidence does not attribute success to the recurrent state, effect score, deferral, or invalidation components. M12 should consequently be described as an implemented composite whose performance hypothesis remains unsupported, not as a validated improvement.


\begin{algorithm}[H]\caption{Ordered guarded selection with broker-admitted candidates}\label{alg:guarded}
\begin{algorithmic}[1]\Require Observation $o_t$, candidates $A_t$, admission predicate $G$
\State Refresh stale observation; acknowledge duplicate events
\State Resolve missing dependency and cancellation preconditions
\For{each operation guard $R_q$ in declared priority order}
  \If{$R_q(o_t)$ and an admitted candidate has operation $q$}
    \State Return the first such candidate
  \EndIf
\EndFor
\State Return typed stop with reason no-supported-rule
\end{algorithmic}\end{algorithm}
\begin{algorithm}[H]\caption{Asynchronous planning with binding validation and fallback accounting}\label{alg:async}
\begin{algorithmic}[1]\Require Observation $o_t$, binding $b_t$, fallback policy $F$
\If{no request is pending and the decision is eligible}
  \State Submit planner request bound to $H(\operatorname{canonical}(b_t))$
\EndIf
\State Continue with $F(o_t)$ while the request is pending
\If{planner result is ready}
  \If{binding digest matches and candidate validates}
    \State Return planner proposal for fresh broker admission
  \Else
    \State Discard stale result and record unused computation
  \EndIf
\EndIf
\State Return fallback; account outstanding work at episode close
\end{algorithmic}\end{algorithm}
\begin{algorithm}[H]\caption{Resource-versioned metacontrol and stale-response handling}\label{alg:m12}
\begin{algorithmic}[1]\Require Observation $o_t$, history $h_t$, candidates $A_t$, budget $B_t$
\State Refresh, acknowledge duplicates, and resolve dependencies
\State Score admitted actions using recurrent preference and effect estimate
\State Update prediction residual $e_t$ from the next observable state
\If{benefit gate, surprise, or explicit uncertainty requests deliberation}
  \If{request budget remains and $B_t$ exceeds minimum}
    \State Submit request bound to current resource revision and candidates
  \EndIf
\EndIf
\If{pending answer is ready and its resource binding remains current}
  \State Validate candidate and submit it to broker
\Else
  \State Discard stale or expired response; choose highest-scoring admitted action
\EndIf
\end{algorithmic}\end{algorithm}

\section{Experimental protocol}
The core matrix has twelve methods, twenty families, six variants, and ten seeds per method-family-variant cell: $12\times20\times6\times10=14{,}400$ episodes. Families cover file routing, builds, dependencies, HTTP retries, process control, conflicts, planner use, untrusted text, long-job cancellation, and crash recovery. Canonical traces bind method, case, variant, seed, decisions, metrics, postcondition, and source lineage. The derived index includes source digests. Figure~\ref{fig:seed-variation} reports the adequacy result: only 71 of 1,440 method-case-variant cells vary across seeds, so these repetitions chiefly test deterministic replay rather than independent environment sampling.
\begin{equation}\label{eq:scheduled-core}N_{\mathrm{core}}=|\mathcal M|\,|\mathcal C|\,|\mathcal V|\,|\mathcal S|=12\cdot20\cdot6\cdot10=14{,}400.\end{equation}
\begin{figure}[H]\centering\includegraphics[width=.82\linewidth]{seed-variation.pdf}\caption{Case-by-variant cells with invariant and seed-varying binary outcomes show why repeated episodes do not imply independent evidence.}\label{fig:seed-variation}\end{figure}

Failures stay in denominators. Success is the count of true independent postconditions divided by scheduled episodes. Decision latency is measured per controller decision; total wall time includes additional work and is reported separately. Paired comparisons match case, variant, and seed. Bootstrap intervals over seed groups can be degenerate when outcomes do not vary, and cannot establish generalization from repeated deterministic episodes. Figure~\ref{fig:core-success} shows the ceiling effect directly: several methods solve every core episode, including a deterministic baseline.

Structural-transfer compositions and semantic-transfer analyses are separate populations, not part of the core count. Four ablations remove memory, effect information, deferral, or invalidation, each over 1,200 matched core episodes. Their outcomes are paired against full M12. Planner cost, policy latency, and task completion are distinct. The frozen protocol and source hashes identify these populations; they do not turn seed relabeling into additional environmental variation. Future experiments must split by task family before fitting and keep final labels inaccessible until evaluation \cite{efron1993,demsar2006,nadeau2003,toolsandbox,swebench,taubench}. Training, calibration, threshold selection, and final evaluation use separate records where available; no held-out outcome should tune a request threshold. For asynchronous planners, request count and applied-decision count are both recorded, because a request may complete too late or be discarded after invalidation. Transfer lane indices preserve the split boundary, and report digests identify the exact artifact snapshot. This protocol supports replay of the archived evaluation but its low seed variability limits inferential resolution. For transfer and routing, the index digest and lineage signature are checked independently before the outcome table is generated. The resulting comparisons describe the declared held-out population, not unrestricted tool use or deployed user workloads. The counts are reviewed in the context of family overlap and the saturated fixed-policy baseline before any generalization language is used.
